%% ============================================================
%%  main.tex  —  ACM TOPLAS submission
%%  lean-pl-verify: Unified Monadic Verification of Rust and TypeScript
%%
%%  Build:  pdflatex main && bibtex main && pdflatex main && pdflatex main
%% ============================================================

\documentclass[acmsmall,screen=true,nonacm]{acmart}

%% ---- packages -----------------------------------------------
\usepackage{booktabs}
\usepackage{microtype}
\usepackage[scaled=0.85]{beramono}  % monospace font with clear I/O/1/0 distinction
\usepackage{amsmath}           % amssymb already loaded by acmart
\usepackage{xcolor}
\usepackage{listings}
\usepackage{url}
\usepackage{enumitem}
\usepackage{xspace}
\usepackage{graphicx}
\usepackage{multirow}
\usepackage{subcaption}
% stmaryrd omitted to avoid \Bbbk conflict with acmart

%% ---- citation style -----------------------------------------
\citestyle{acmauthoryear}
\setcitestyle{nosort}

%% ---- Lean 4 listing style -----------------------------------
\definecolor{lean-keyword}{rgb}{0.10, 0.10, 0.55}
\definecolor{lean-comment}{rgb}{0.30, 0.50, 0.30}
\definecolor{lean-string}{rgb}{0.55, 0.10, 0.10}
\definecolor{lean-type}{rgb}{0.05, 0.35, 0.55}
\definecolor{lean-tactic}{rgb}{0.50, 0.10, 0.50}
\definecolor{codebg}{rgb}{0.97, 0.97, 0.97}

\lstdefinelanguage{Lean4}{
  morekeywords={def,theorem,lemma,example,let,have,show,by,where,import,
    namespace,end,open,section,variable,instance,class,structure,inductive,
    match,with,fun,if,then,else,return,do,for,while,in,at,and,or,not,
    intro,exact,apply,simp,rfl,omega,linarith,norm_num,ring,tauto,
    decide,native_decide,refine,rw,rewrite,cases,induction,contradiction,
    assumption,trivial,constructor,ext,calc,conv,suffices,obtain,
    private,protected,noncomputable,abbrev,attribute,macro,syntax,
    macro_rules,notation,infixl,infixr,prefix},
  sensitive=true,
  morecomment=[l]{--},
  morecomment=[s]{/-}{-/},
  morestring=[b]",
  keywordstyle=\color{lean-keyword}\bfseries,
  commentstyle=\color{lean-comment}\itshape,
  stringstyle=\color{lean-string},
  basicstyle=\ttfamily\small,
  backgroundcolor=\color{codebg},
  frame=single,
  framerule=0.4pt,
  rulecolor=\color{gray!50},
  xleftmargin=6pt,
  xrightmargin=6pt,
  breaklines=true,
  showstringspaces=false,
  tabsize=2,
  captionpos=b,
  numbers=left,
  numberstyle=\tiny\color{gray},
  numbersep=5pt,
  literate=
    {≤}{{\ensuremath{\leq}}}1
    {≥}{{\ensuremath{\geq}}}1
    {≠}{{\ensuremath{\neq}}}1
    {←}{{\ensuremath{\leftarrow}}}1
    {→}{{\ensuremath{\rightarrow}}}1
    {↑}{{\ensuremath{\uparrow}}}1
    {↓}{{\ensuremath{\downarrow}}}1
    {∧}{{\ensuremath{\wedge}}}1
    {∨}{{\ensuremath{\vee}}}1
    {∀}{{\ensuremath{\forall}}}1
    {∃}{{\ensuremath{\exists}}}1
    {α}{{\ensuremath{\alpha}}}1
    {β}{{\ensuremath{\beta}}}1
    {σ}{{\ensuremath{\sigma}}}1
    {⟨}{{\ensuremath{\langle}}}1
    {⟩}{{\ensuremath{\rangle}}}1
    {·}{{\ensuremath{\cdot}}}1
    {↦}{{\ensuremath{\mapsto}}}1
    {⊢}{{\ensuremath{\vdash}}}1
    {∈}{{\ensuremath{\in}}}1
    {ℕ}{{\ensuremath{\mathbb{N}}}}1
    {⟹}{{\ensuremath{\Longrightarrow}}}1,
}

\lstset{language=Lean4,
  literate=
    % --- Unicode symbols (must come before ASCII digraphs) ---
    {≤}{{\ensuremath{\leq}}}1
    {≥}{{\ensuremath{\geq}}}1
    {≠}{{\ensuremath{\neq}}}1
    {←}{{\ensuremath{\leftarrow}}}1
    {→}{{\ensuremath{\rightarrow}}}1
    {↑}{{\ensuremath{\uparrow}}}1
    {↓}{{\ensuremath{\downarrow}}}1
    {∧}{{\ensuremath{\wedge}}}1
    {∨}{{\ensuremath{\vee}}}1
    {∀}{{\ensuremath{\forall}}}1
    {∃}{{\ensuremath{\exists}}}1
    {α}{{\ensuremath{\alpha}}}1
    {β}{{\ensuremath{\beta}}}1
    {σ}{{\ensuremath{\sigma}}}1
    {⟨}{{\ensuremath{\langle}}}1
    {⟩}{{\ensuremath{\rangle}}}1
    {·}{{\ensuremath{\cdot}}}1
    {↦}{{\ensuremath{\mapsto}}}1
    {⊢}{{\ensuremath{\vdash}}}1
    {∈}{{\ensuremath{\in}}}1
    {ℕ}{{\ensuremath{\mathbb{N}}}}1
    {⟹}{{\ensuremath{\Longrightarrow}}}1
    % --- ASCII digraphs ---
    {<=}{{\ensuremath{\leq}}}1
    {>=}{{\ensuremath{\geq}}}1
    {/=}{{\ensuremath{\neq}}}1
    {->}{{\ensuremath{\to}}}2
    {<-}{{\ensuremath{\leftarrow}}}2
    {:=}{{\ensuremath{:=}}}2
    {|=}{{\ensuremath{\models}}}2
    {sigma}{{\ensuremath{\sigma}}}5
    {alpha}{{\ensuremath{\alpha}}}5
    {beta}{{\ensuremath{\beta}}}4
}

%% ---- TypeScript listing style ------------------------------
\lstdefinelanguage{TypeScript}{
  morekeywords={function,return,let,const,var,if,else,while,for,break,continue,
    true,false,null,undefined,typeof,instanceof,new,this,class,extends,
    import,export,from,async,await,number,boolean,string,void},
  sensitive=true,
  morecomment=[l]{//},
  morecomment=[s]{/*}{*/},
  morestring=[b]",
  morestring=[b]',
  keywordstyle=\color{lean-keyword}\bfseries,
  commentstyle=\color{lean-comment}\itshape,
  stringstyle=\color{lean-string},
  basicstyle=\ttfamily\small,
  backgroundcolor=\color{codebg},
  frame=single,
  framerule=0.4pt,
  rulecolor=\color{gray!50},
  xleftmargin=6pt,
  xrightmargin=6pt,
  breaklines=true,
  showstringspaces=false,
  tabsize=2,
}

%% ---- macros -------------------------------------------------
\newcommand{\RustM}{\texttt{RustM}\xspace}
\newcommand{\ProgramSpec}{\texttt{ProgramSpec}\xspace}
\newcommand{\satisfies}{\mathbin{\models}}
\newcommand{\evalF}[1]{\texttt{evalStmtFuel}\,#1}
\newcommand{\evalFun}{\texttt{evalFun}\xspace}
\newcommand{\LLBCStmt}{\texttt{LLBCStmt}\xspace}
\newcommand{\Locals}{\texttt{Locals}\xspace}
\newcommand{\sorry}{\texttt{sorry}\xspace}
\newcommand{\rfl}{\texttt{rfl}\xspace}
\newcommand{\simp}{\texttt{simp}\xspace}
\newcommand{\omegaT}{\texttt{omega}\xspace}
\newcommand{\lean}[1]{\lstinline[language=Lean4]{#1}}
\newcommand{\lakeCmd}[1]{\texttt{\small #1}}
\newcommand{\minI}{\texttt{minI32}\xspace}
\newcommand{\maxI}{\texttt{maxI32}\xspace}
\newcommand{\FibOv}{\texttt{FibOv}\xspace}
\newcommand{\SumToOv}{\texttt{SumToOv}\xspace}
\newcommand{\FactOv}{\texttt{FactOv}\xspace}
\newcommand{\natfib}{\texttt{Nat.fib}\xspace}
\newcommand{\ie}{\textit{i.e.},\xspace}
\newcommand{\eg}{\textit{e.g.},\xspace}
\newcommand{\vs}{\textit{vs.}\xspace}
\newcommand{\etal}{\textit{et al.}\xspace}

%% ---- theorem environments -----------------------------------
\newtheorem{definition}{Definition}[section]
\newtheorem{theorem}{Theorem}[section]
\newtheorem{lemma}[theorem]{Lemma}
\newtheorem{corollary}[theorem]{Corollary}
\newtheorem{remark}{Remark}[section]

%% ============================================================
%%  Paper metadata
%% ============================================================

\title{A Unified Monadic Framework for Cross-Language Formal Verification
       in Lean~4: Deep LLBC Embedding, Shared Specifications,
       and Proof Automation}

\author{Anand Kumar Keshavan}
\orcid{0009-0007-8541-5203}
\affiliation{%
  \institution{Independent Researcher}
  \city{Pune}
  \country{India}
}
\email{[author email redacted for review]}

%% ---- CCS concepts -------------------------------------------
\begin{CCSXML}
<ccs2012>
 <concept>
  <concept_id>10003752.10003790.10003795</concept_id>
  <concept_desc>Theory of computation~Proof theory</concept_desc>
  <concept_significance>500</concept_significance>
 </concept>
 <concept>
  <concept_id>10003752.10003790.10011740</concept_id>
  <concept_desc>Theory of computation~Type theory</concept_desc>
  <concept_significance>300</concept_significance>
 </concept>
 <concept>
  <concept_id>10011007.10011006.10011008.10011024.10011027</concept_id>
  <concept_desc>Software and its engineering~Formal methods</concept_desc>
  <concept_significance>500</concept_significance>
 </concept>
 <concept>
  <concept_id>10011007.10011006.10011008.10011024</concept_id>
  <concept_desc>Software and its engineering~Software verification</concept_desc>
  <concept_significance>500</concept_significance>
 </concept>
 <concept>
  <concept_id>10011007.10011006.10011050.10011017</concept_id>
  <concept_desc>Software and its engineering~Domain specific languages</concept_desc>
  <concept_significance>300</concept_significance>
 </concept>
</ccs2012>
\end{CCSXML}

\ccsdesc[500]{Theory of computation~Proof theory}
\ccsdesc[300]{Theory of computation~Type theory}
\ccsdesc[500]{Software and its engineering~Formal methods}
\ccsdesc[500]{Software and its engineering~Software verification}
\ccsdesc[300]{Software and its engineering~Domain specific languages}

\keywords{Lean 4, formal verification, Rust, TypeScript, monadic semantics,
          LLBC, loop invariants, proof automation, deep embedding,
          ProgramSpec, cross-language verification}

%% ---- abstract -----------------------------------------------
\begin{abstract}
Formal verification of systems software typically proceeds language by language:
a tool is built for Rust, another for C, another for TypeScript. Each brings
its own trusted translation pipeline, its own specification logic, and its own
proof methodology. The cost of entry — and the cost of porting properties across
language boundaries — is high.

This paper describes \emph{lean-pl-verify}, a Lean~4 library that takes a
different approach. Rather than building a language-specific verifier, we
define a shared execution monad (\RustM) and a shared specification language
(\ProgramSpec) that are independent of any particular source language. Programs
written in Rust and TypeScript are independently translated into this common
semantic target via deep embeddings: a full inductive encoding of the LLBC
intermediate representation for Rust, and a parallel inductive encoding of a
TypeScript AST. Once in this shared form, all verification tools — specification
combinators, loop invariant patterns, proof automation macros — apply uniformly.

The library contains \textbf{258 kernel-checked theorems} (0 \sorry) proved in Lean~4,
with no trusted code generator in the trusted computing base. Every theorem is
verified by the Lean kernel; the single command \lakeCmd{lake build Theorems}
reproduces all results.

Our chief novel findings are:
(1)~a fuel-parameterized deep LLBC interpreter definitionally transparent to
the Lean kernel, enabling \rfl proofs for pure functions;
(2)~a relational big-step semantics with a \emph{full adequacy} proof (18 theorems,
0 \sorry) establishing both soundness and completeness of the interpreter
with respect to a language-independent denotational specification;
(3)~an end-to-end Charon pipeline (\texttt{charon2lean.py}, 463 lines) with
57 theorems over 18 generated definitions, plus 5 theorems for
\texttt{Integer::is\_even<u32>} extracted from the published
\texttt{num-integer} v0.1.45 crate (${\approx}$45M downloads on crates.io);
(4)~four proof-automation macros that reduce routine proof overhead by $12{\times}$;
(5)~39 TypeScript theorems — including a verified \texttt{while} loop and a
bug-detection case study exposing two off-by-one and wrong-guard bugs — using
the same monad and specification as Rust; and
(6)~six cross-language unification theorems showing that identical \ProgramSpec
specifications are satisfied by both Rust and TypeScript implementations.

\bigskip
\noindent\textit{Artifact statement.} The complete source is available as
a supplementary artifact. Build prerequisites: Lean~4.30.0-rc2 (via \texttt{elan})
and Mathlib4. Verification: \lakeCmd{lake build Theorems}.
\end{abstract}

%% ============================================================
\begin{document}
\maketitle

%% ---- AI disclosure ------------------------------------------
\begin{acks}
The research design, all theorem statements, proof strategies, architectural
decisions, and scientific claims are the author's own.

\textbf{AI tool disclosure.} Claude Code (Anthropic, powered by Claude~Sonnet~4.6),
an AI coding agent, was used extensively throughout this project in the following
roles: (i)~generating and iterating on Lean~4 proof scripts, tactic sequences,
and simp sets, with all outputs verified by \lakeCmd{lake build} before acceptance;
(ii)~writing the 463-line \texttt{charon2lean.py} translator pipeline almost
entirely, with the author directing the overall approach and resolving
installation blockers; (iii)~scaffolding proof structures (\eg the adequacy
induction skeleton) once the author specified the top-level theorem statement;
(iv)~generating repetitive proof instances across multiple functions from an
author-supplied template; and (v)~drafting and revising sections of this paper
in response to author direction, with all scientific statements reviewed and
corrected by the author.

No theorem appears in this paper without independent kernel verification.
The Lean kernel, not the AI agent, is the trust anchor: \lakeCmd{lake build Theorems}
is the reproducibility certificate for all 258 theorems (0 sorry).
The author takes full responsibility for all scientific claims.
\end{acks}

%% ============================================================
%%  §1  Introduction
%% ============================================================

\section{Introduction}
\label{sec:intro}

Modern software stacks mix languages. A backend service might be written in
Rust for performance-critical subsystems, TypeScript for the API layer, and
Python for tooling. When the goal is formal verification, this heterogeneity
poses a quiet but persistent problem: each existing verification framework is
coupled to one language. Aeneas~\cite{ho2022aeneas} targets safe Rust through
functional translation. Verus~\cite{verus2023} provides deductive verification
for Rust via SMT-backed contracts. Creusot~\cite{creusot2022} translates Rust
into Why3 proof obligations. None of these tools say anything about TypeScript,
and none of their specifications are portable across language boundaries.

The result is that proving a property about a Rust function and then extending
that proof to a semantically equivalent TypeScript function requires starting
over: a different tool, a different specification language, a different proof
style. Cross-language correctness — the claim that two implementations of the
same function, in different languages, satisfy the same functional
specification — is effectively unprovable with the current toolkit.

\subsection{The Key Insight}

Rust and TypeScript differ structurally. Rust's ownership system enforces
linear resource discipline at compile time; TypeScript runs on a garbage-collected
runtime with dynamic typing and implicit coercions. Their compilers produce
radically different intermediate representations. Yet at the level of pure
functional semantics, both execute sequences of operations that either
succeed (producing a value) or fail (raising an exception or panicking). That
shared structure is precisely what a state-error monad captures.

If programs from both languages are translated into a single monadic target —
a monad that models local state and runtime failure — then a specification
written over that monad describes both programs without modification. This
observation drives the entire design of \emph{lean-pl-verify}.

\subsection{Our Approach}

We define two components that are intentionally language-agnostic:

\begin{itemize}[leftmargin=1.5em]
  \item \textbf{\RustM}, the execution monad:
        $\RustM\;\sigma\;\alpha \;=\; \texttt{StateT}\;\sigma\;(\texttt{Except}\;\texttt{PanicReason})\;\alpha$.
        This represents a computation that may read and update a state $\sigma$,
        may return a value of type $\alpha$, and may instead signal a panic.

  \item \textbf{\ProgramSpec}, the specification language: an inductive type with
        five constructors (\texttt{pureOutput}, \texttt{nocrash}, \texttt{both},
        \texttt{withFuel}, \texttt{precond}) and a satisfaction relation
        $m \satisfies P$ defined by recursion on $P$.
\end{itemize}

Programs from Rust and TypeScript are translated into \RustM via separate deep
embeddings. A \emph{deep embedding} means that the source language's syntax is
represented as an inductive Lean~4 datatype, and execution is defined as a
definitional function over that datatype. The Lean kernel can therefore evaluate
closed programs directly, which is what enables \rfl-style proofs for pure
functions. No external trusted component sits between the source program and
the Lean kernel.

For Rust we embed a substantial fragment of LLBC (Low-Level Borrow
Calculus)~\cite{charon2024}: the AST covers assignments, conditionals,
sequential composition, loops with \texttt{break}/\texttt{continue},
and explicit function calls. For TypeScript we embed a parallel AST covering
expression evaluation, binary operations, conditionals, and function calls.
Both ASTs use the same \Locals representation (a list of \texttt{Value}s),
the same \RustM monad, and the same \ProgramSpec specification language.

\subsection{Contributions}

This paper makes the following contributions.

\begin{enumerate}[label=\textbf{C\arabic*.},leftmargin=2em]
  \item \textbf{RustM monad} (\S\ref{sec:monad}). A kernel-transparent
        state-error monad with a \texttt{PanicReason} inductive type covering
        arithmetic overflow, division by zero, explicit panics, and index
        errors. Two proved monad laws (\texttt{bind\_ok}, \texttt{pure\_ok})
        and a forward-backward model for mutable references.

  \item \textbf{Deep LLBC embedding} (\S\ref{sec:llbc}). An inductive Lean~4
        encoding of the LLBC AST with a fuel-parameterized interpreter
        \texttt{evalStmtFuel} that is definitionally transparent to the kernel.
        The interpreter handles the full LLBC fragment used by Charon~\cite{charon2024}
        for safe Rust: loops, break, conditionals, multi-function environments.

  \item \textbf{ProgramSpec} (\S\ref{sec:programspec}). A language-agnostic
        specification combinator library with 17 kernel-checked theorems
        (\texttt{sat\_pure}, \texttt{sat\_bind}, \texttt{sat\_both\_intro},
        \texttt{sat\_pureOutput\_mono}, and~13 others). The library makes no
        reference to LLBC or TypeScript ASTs.

  \item \textbf{Forward-backward ownership model} (\S\ref{sec:monad}). A pure
        functional encoding of \texttt{\&mut T} via \texttt{FBPair} with four
        identity laws proved in Lean, matching the formal semantics of
        Aeneas~\cite{ho2022aeneas} without requiring separation logic.

  \item \textbf{Loop invariant methodology} (\S\ref{sec:rust-proofs}). A
        reusable inductive invariant pattern demonstrated on three loop
        functions: \texttt{sumTo} (Gauss formula), \texttt{fact} (factorial),
        and \texttt{fib} (Fibonacci, \S\ref{sec:fib}). The pattern involves an
        explicit overflow predicate, a parametric loop-correct theorem proved
        by induction on remaining iterations, and a safety-implies-overflow
        bridge lemma.

  \item \textbf{Proof automation} (\S\ref{sec:automation}). Four Lean~4 tactic
        macros (\texttt{llbc\_verify}, \texttt{llbc\_verify\_prop},
        \texttt{llbc\_verify\_cond}, \texttt{llbc\_verify\_loop}) that reduce
        typical 12-line LLBC proofs to single-line calls. Measured compression:
        12$\times$.

  \item \textbf{Relational semantics and full adequacy} (\S\ref{sec:semantics}).
        A relational big-step semantics (\texttt{EvalStmt}, \texttt{EvalFun},
        and four sub-relations) provides a language-independent ground truth for
        what LLBC programs mean. A \emph{full adequacy} theorem (18 kernel-checked
        proofs, 0 \sorry) establishes both soundness (every interpreter output is
        relationally derivable) and completeness (every relational derivation is
        witnessed by some fuel bound), closing the gap between the fuel-based
        evaluator and a declarative specification.

  \item \textbf{Charon end-to-end pipeline and real-crate extraction}
        (\S\ref{sec:eval-charon}--\ref{sec:eval-numinteger}).
        A 463-line Python translator (\texttt{charon2lean.py}) automatically
        converts Charon LLBC JSON to \texttt{LLBCFunDef} Lean terms.
        57 theorems prove the Charon-generated definitions correct across
        18 functions, including \texttt{pow} and \texttt{gcd}.
        Additionally, 5 theorems verify \texttt{Integer::is\_even<u32>}
        extracted from \texttt{num-integer} v0.1.45 (${\approx}$45M downloads),
        a published crate on \url{crates.io}.
        The pipeline: \texttt{verified\_fns.rs} $\to$ Charon $\to$ LLBC JSON $\to$
        \texttt{charon2lean.py} $\to$ \texttt{CharonDefs.lean} $\to$ proofs.

  \item \textbf{TypeScript extension and unification} (\S\ref{sec:typescript}).
        33 theorems for a TypeScript AST interpreter (including a verified
        \texttt{while} loop) using the same \RustM monad and \ProgramSpec
        specification. Six cross-language unification theorems — covering
        \texttt{max}, \texttt{min}, \texttt{mul}, \texttt{add}, \texttt{neg},
        and the looping \texttt{sum\_to} — formally assert that Rust and
        TypeScript implementations satisfy identical specifications.

  \item \textbf{Reproducible artifact} (\S\ref{sec:evaluation}). \textbf{258} kernel-checked
        theorems in 5\,361 lines of Lean~4 code, \textbf{0 sorry}. All theorems are verified
        by a single command: \lakeCmd{lake build Theorems}. The artifact is
        independently reproducible given Lean~4.30.0-rc2 and Mathlib4.
\end{enumerate}

\subsection{What This Work Is Not}

Honesty about scope matters. This work does not verify unsafe Rust: pointer
arithmetic, raw memory, and concurrency are outside the model.
The verification covers functional correctness properties (\eg correctness of
Fibonacci), not security properties or liveness.
The Charon pipeline handles the 16 functions in the artifact; scaling to large
real-world crates is future work.
All 258 theorems are fully kernel-checked: \textbf{0 sorry}.

\subsection{Why Deep Embedding?}

The prevailing approach in tools like Aeneas uses a \emph{shallow} embedding:
Rust functions are compiled to native Lean~4 terms, bypassing an explicit
interpreter. This produces cleaner generated code and shorter proofs for pure
functions. The cost is a trusted code generator: correctness of the generated
Lean depends on correctness of the compiler, which is not itself formally
verified.

Our deep embedding inverts this tradeoff. The interpreter is a Lean~4 \texttt{def},
so every computation step is visible to the Lean kernel. No external component
is trusted. The cost is that proof terms are larger, and loop invariant proofs
require more manual effort. We believe this tradeoff is the right one for a
library whose primary purpose is to establish the \emph{foundations} of
cross-language verification: when the question is ``what do these programs
actually mean?'', having the answer be ``exactly what the Lean kernel says''
is a valuable guarantee.

\subsection{Paper Organisation}

Section~\ref{sec:background} reviews background and related work.
Section~\ref{sec:monad} introduces \RustM and the forward-backward model.
Section~\ref{sec:llbc} presents the LLBC embedding and interpreter.
Section~\ref{sec:programspec} defines \ProgramSpec.
Section~\ref{sec:rust-proofs} walks through the Rust verification results,
including the loop invariant methodology.
Section~\ref{sec:automation} describes the proof automation macros.
Section~\ref{sec:typescript} covers the TypeScript extension and
cross-language unification.
Section~\ref{sec:fib} presents the Fibonacci case study.
Section~\ref{sec:semantics} presents the relational semantics and full adequacy proof.
Section~\ref{sec:evaluation} gives the evaluation.
Section~\ref{sec:agentic} reflects on the agentic development process.
Section~\ref{sec:related} discusses related work.
Section~\ref{sec:conclusion} concludes.

%% ============================================================
%%  §2  Background and Related Work
%% ============================================================

\section{Background and Related Work}
\label{sec:background}

\subsection{Lean 4 as a Verification Platform}

Lean~4~\cite{lean4} is simultaneously a dependently typed programming language
and an interactive theorem prover. The two roles share a common kernel: a small,
trusted type checker that accepts or rejects proof terms. This kernel is the
trust anchor. A theorem is established if and only if the kernel accepts a
proof term of the correct type.

For our purposes, two properties of Lean~4 are critical. First, \emph{kernel
transparency}: definitions marked \texttt{def} (as opposed to \texttt{opaque}
or \texttt{@[extern]}) are unfolded by the kernel during type-checking. This
means that a closed expression like \texttt{evalFun [] return42Fun 5 []}
can be fully reduced by the kernel to its normal form, making proofs by
\texttt{rfl} possible. Second, \emph{definitional equality}: the kernel
checks proof terms by normalisation, so two terms that reduce to the same
normal form are accepted as equal without requiring explicit rewrite steps.

Lean~4's tactic system is implemented in Lean~4 itself, which makes it
extensible. Custom tactics and macros can be written as Lean~4 programs,
compiled, and used immediately within the same file. This facility is
central to the proof automation described in \S\ref{sec:automation}.

\textbf{Mathlib.} The Mathlib library~\cite{mathlib2020} provides an extensive
collection of mathematical results for Lean~4, including the Fibonacci
sequence (\texttt{Nat.fib} and its properties), linear arithmetic tactics
(\texttt{omega}), and algebraic simplification (\texttt{ring}, \texttt{norm\_num}).
Our project imports \texttt{Mathlib.Tactic} for access to these tools;
the Fibonacci case study (\S\ref{sec:fib}) uses
\texttt{Nat.fib\_add\_two} and \texttt{Nat.fib\_mono} from Mathlib.

\subsection{Rust's Ownership Model and LLBC}

Rust enforces an ``aliasing XOR mutability'' discipline via its type
system~\cite{rustOwnership2014}. At any point, a memory location is either
accessible through multiple immutable references, or through exactly one mutable
reference, but not both. This invariant is checked statically by the borrow
checker and eliminates the data races and use-after-free bugs that afflict C and
C++ programs.

The compiler pipeline proceeds: Rust source $\to$ HIR (high-level IR)
$\to$ MIR (mid-level IR) $\to$ LLVM IR. MIR makes borrow operations, moves,
and drops explicit, but its control-flow graph is unstructured (arbitrary jumps,
implicit panic handlers). The Charon framework~\cite{charon2024} consumes MIR
and produces two higher-level IRs: ULLBC (Unstructured LLBC), which simplifies
pattern matching and constants, and LLBC (Low-Level Borrow Calculus), which
reconstructs structured control flow — loops, matches, sequential branches.
LLBC is the target of our deep embedding.

The key property of LLBC for our purposes is that loops are represented as
a \emph{single} construct: \texttt{loop body} executes \texttt{body}
repeatedly until a \texttt{break} statement is encountered. This maps directly
to the recursive case in our fuel-parameterised interpreter. Conditionals are
represented as \texttt{ite cond trueBranch falseBranch}. Assignments are
\texttt{assign dest rvalue cont}, where the continuation \texttt{cont} is the
next statement.

\subsection{The Forward-Backward Paradigm for \texttt{\&mut T}}

Mutable references in Rust present a semantic challenge for functional models.
When a function returns a \texttt{\&'a mut T}, the caller may subsequently
mutate through that reference, and those mutations must propagate back to
the original storage location. Representing this naively requires a heap model
or pointer arithmetic.

The Aeneas translation~\cite{ho2022aeneas} resolves this with the
\emph{forward-backward} paradigm. A function that operates on mutable borrows
is split into: a \emph{forward function} (modelling execution up to the return
point) and a \emph{backward function} (propagating borrow updates back to
inputs when lifetimes end). Both are pure. We implement the same paradigm
via the \texttt{FBPair} type in \S\ref{sec:monad}.

\subsection{Comparison with Existing Rust Verification Systems}

Table~\ref{tab:comparison} summarises the relationship between lean-pl-verify
and the main competing tools. Several distinctions merit comment.

\begin{table*}[t]
\centering
\caption{Comparison of Rust verification approaches.
         ``TCB'' = trusted computing base; ``Lang.'' = source language(s) supported;
         ``Loop'' = support for unbounded loop verification.}
\label{tab:comparison}
\small
\begin{tabular}{@{}llllll@{}}
\toprule
\textbf{Tool} & \textbf{Approach} & \textbf{TCB} & \textbf{Lang.} & \textbf{Loop} & \textbf{Notes} \\
\midrule
Aeneas~\cite{ho2022aeneas}   & Shallow monadic translation & Charon + code generator & Safe Rust & Yes & Loop summaries \\
Verus~\cite{verus2023}        & SMT-backed contracts & Z3 + verifier & Rust (incl.\ unsafe) & Yes & SMT cliff on complex goals \\
Creusot~\cite{creusot2022}   & Why3 deductive & Why3 + SMT & Safe Rust & Yes & Intrinsic specification \\
Kani~\cite{kani2022}         & Bounded model checking & CBMC/SAT & Rust (incl.\ unsafe) & Bounded & No unbounded loops \\
RustBelt~\cite{rustbelt2018} & Semantic (Iris/Coq) & Coq + Iris & Unsafe Rust & Yes & High manual cost \\
\midrule
\textbf{This work}            & Deep embedding (Lean 4) & Lean kernel only & Rust + TypeScript & Yes & Cross-language unification \\
\bottomrule
\end{tabular}
\end{table*}

\textbf{Aeneas} is the closest technical relative. Both Aeneas and our work
translate safe Rust into Lean~4 and prove functional correctness properties.
The key difference is the embedding strategy. Aeneas generates native Lean
functions via a trusted translation pipeline; correctness of the generated code
depends on correctness of the Charon+Aeneas toolchain. Our deep embedding
makes no such assumption: the interpreter is a Lean~4 \texttt{def}, and its
correctness is definitionally visible to the kernel. A second difference is
scope: Aeneas does not target TypeScript, and cross-language unification is
not a goal. A third difference is proof style: Aeneas encourages users to
write proofs about generated Lean functions, whereas we write proofs directly
about the AST.

\textbf{Verus} offers the strongest industrial story among the compared tools,
with real-world verification of distributed systems components. Its reliance on
Z3 introduces the possibility of false negatives (SMT timeouts) and places the
SMT solver in the TCB. Verus does not support TypeScript.

\textbf{F$^\star$}~\cite{fstar2016} deserves mention: it uses Dijkstra monads
to derive verification condition generators from effectful programs, placing
it philosophically close to our use of \RustM. The difference is that F$^\star$
targets its own language; applying it to Rust would require a different extraction
pipeline.

\textbf{LOOM}~\cite{loom2026} is a recent framework for foundational multi-modal
verification using Dijkstra monads in Lean~4. It shares our interest in
language-agnostic specification but does not target Rust or TypeScript programs
via deep AST embeddings; it operates on shallow monadic terms.

\subsection{Deep \vs Shallow Embeddings}

The choice between deep and shallow embeddings in mechanised semantics is
well-studied~\cite{coqArt}. A shallow embedding maps source constructs directly
to host language constructs; proofs about the source are proofs about the host
language, benefiting from all of the host's proof support. A deep embedding
represents source syntax explicitly; proofs require explicitly reasoning about
the interpreter, but the syntax is a first-class object that can be inspected
and transformed.

For our use case, the deep embedding brings one decisive advantage: the Lean
kernel can evaluate the interpreter on closed terms. A ground computation
like \texttt{evalFun [] fibFun 50 [.int 10]} reduces by kernel reduction
to a concrete result, and \texttt{rfl} proves the correctness lemma. With
a shallow embedding, this would require rewriting along the compiled function
body, which is generally less direct.

\subsection{Note on AI-Assisted Development}

This paper was developed with substantial assistance from Claude Code
(Anthropic), an AI coding agent. The research design, theorem statements,
proof strategies, and scientific conclusions are the author's own. The AI
agent's roles included: generating Lean~4 proof scripts (all verified by
\lakeCmd{lake build}), writing the Charon translation pipeline, scaffolding
proof structures from author-specified statements, generating repetitive proof
instances, and drafting paper text under author direction.

Readers should be aware that a significant portion of the prose and code was
AI-generated and subsequently reviewed and corrected by the author. Variations
in writing style between sections reflect the iterative human-AI editing
process. The complete AI disclosure appears in the Acknowledgements.

All 258 theorems are kernel-checked (0 sorry); no theorem was accepted without
\lakeCmd{lake build} verification. The Lean kernel is the sole trust anchor.

%% ============================================================
%%  §3  The Execution Monad
%% ============================================================

\section{The Execution Monad}
\label{sec:monad}

The semantic foundation of lean-pl-verify is a state-error monad that captures
the observable behaviour of safe Rust and TypeScript programs: they compute
with local state, and they may panic. This section presents the monad
definition, its basic properties, and the forward-backward model for mutable
references.

\subsection{Panic Reasons and Integer Bounds}
\label{sec:types}

Rust programs can panic at runtime for several distinct reasons. We enumerate
them as an inductive type:

\begin{lstlisting}[caption={PanicReason inductive type.},label={lst:panic}]
inductive PanicReason where
  | arithmeticOverflow
  | divisionByZero
  | explicitPanic
  | indexOutOfBounds
\end{lstlisting}

This enumeration is not exhaustive with respect to production Rust (stack
overflow and foreign panics are excluded), but it covers every panic that
can arise in the safe Rust fragment we target.

Integer arithmetic is modelled over \texttt{Int} (Lean's unbounded integer
type), with explicit bounds checks at the point of each operation. The bounds
for 32-bit signed integers are:

\begin{lstlisting}[caption={Integer bounds.},label={lst:bounds}]
namespace IntBounds
  def minI32 : Int := -(2^31)     -- -2147483648
  def maxI32 : Int :=  2^31 - 1  --  2147483647
end IntBounds
\end{lstlisting}

Using \texttt{Int} rather than a fixed-width type (\eg Lean's \texttt{Int32})
simplifies the proof theory significantly: the \texttt{omega} and \texttt{linarith}
tactics operate natively on \texttt{Int}, whereas proofs about \texttt{Int32}
would require constant casts and bitwidth reasoning. The trade-off is that
overflow is a proof obligation rather than a type invariant; every arithmetic
theorem must carry explicit preconditions of the form
$\texttt{minI32} \le v \wedge v \le \texttt{maxI32}$.

\subsection{The RustM Monad}

\begin{lstlisting}[caption={RustM monad definition.},label={lst:rustm}]
def RustM (s : Type) (a : Type) :=
  StateT s (Except PanicReason) a
\end{lstlisting}

\texttt{RustM}~$\sigma$~$\alpha$ is a computation that: reads and writes a
state of type $\sigma$ (modelling local variables), may return a value of
type $\alpha$, and may instead panic with a \texttt{PanicReason}. The
\texttt{StateT} transformer and \texttt{Except} type are both defined in
Lean~4's core library as \texttt{def}s, so they are kernel-transparent.

Monad operations are inherited from the transformer stack:

\begin{lstlisting}[caption={Derived monad operations.},label={lst:monadops}]
-- Fail with a panic reason
def rpanic (r : PanicReason) : RustM s a :=
  fun _ => .error r

-- Safe division: panics on zero denominator
def rdiv (a b : Int) : RustM s Int :=
  if b = 0 then rpanic .divisionByZero
  else pure (a / b)

-- Checked addition: panics on overflow
def raddChecked (a b : Int) : RustM s Int :=
  let v := a + b
  if minI32 <= v /\ v <= maxI32 then pure v
  else rpanic .arithmeticOverflow
\end{lstlisting}

Two monad laws are proved as standalone lemmas and used in the
\ProgramSpec reasoning in \S\ref{sec:programspec}:

\begin{lstlisting}[caption={Monad laws for RustM.},label={lst:monadlaws}]
theorem bind_ok {s a b : Type}
    (f : RustM s a) (g : a -> RustM s b)
    (s : s) (a : a) (s' : s)
    (h : f s = .ok (a, s')) :
    (f >>= g) s = g a s' := by
  simp [bind, StateT.bind, h]

theorem pure_ok {s a : Type} (a : a) (s : s) :
    (pure a : RustM s a) s = .ok (a, s) := rfl
\end{lstlisting}

These look trivial, but their role is structural: they let the \ProgramSpec
combinators in \S\ref{sec:programspec} reason about bound computations without
re-expanding monad definitions in every proof.

\subsection{Kernel Transparency}

A critical design constraint is that \RustM must be \emph{kernel-transparent}:
the Lean kernel must be able to reduce expressions of type
\texttt{RustM\;$\sigma$\;$\alpha$} to their normal forms during type-checking.

This is satisfied by definition: \texttt{StateT}, \texttt{Except}, and all
derived operations are \texttt{def}s with no \texttt{@[extern]} or
\texttt{@[opaque]} attributes. Lean's kernel reduction will unfold all of
them when checking a proof term. The practical consequence is that proofs of
the form \texttt{evalFun [] f 30 args at init |= P} can often be closed
by \texttt{rfl} when \texttt{P} is a \texttt{pureOutput} spec and the
evaluation terminates concretely.

This transparency extends to our \emph{Value} type:

\begin{lstlisting}[caption={Value and Locals types.},label={lst:value}]
inductive Value where
  | int   : Int  -> Value
  | bool_ : Bool -> Value
  | unit  : Value

abbrev Locals := List Value
\end{lstlisting}

Using \texttt{List Value} for local variable storage means that all variable
accesses (\texttt{List.getElem?}, \texttt{List.set}) are also kernel-transparent
standard library functions.

\subsection{The Forward-Backward Model for Mutable References}
\label{sec:fbpair}

Safe Rust frequently passes mutable references as function arguments. Consider
the canonical example from the Aeneas literature~\cite{ho2022aeneas}:

\begin{verbatim}
fn choose<'a, T>(b: bool, x: &'a mut T, y: &'a mut T)
    -> &'a mut T {
  if b { x } else { y }
}
\end{verbatim}

The function returns a mutable reference to one of its inputs. The caller
may then mutate through the returned reference; the mutation must propagate
back to \texttt{x} or \texttt{y} depending on which branch was taken.

In a purely functional model, this is encoded using a pair of functions:

\begin{lstlisting}[caption={FBPair: forward-backward representation.},label={lst:fbpair}]
structure FBPair (a b : Type) where
  fwd : a -> b
  bwd : a -> b -> a

-- forward: b=true => fwd selects x; b=false => fwd selects y
-- backward: propagates updated value back to the chosen input
\end{lstlisting}

The \texttt{fwd} function (forward) computes the result of selecting a branch.
The \texttt{bwd} function (backward) takes the original input and the updated
output and returns the updated input — modelling the lifetime-end propagation.

Four identity laws are proved:

\begin{lstlisting}[caption={FBPair identity laws.},label={lst:fblaws}]
theorem fwd_true (x y : a) :
    choose_fwd true x y = x := rfl

theorem fwd_false (x y : a) :
    choose_fwd false x y = y := rfl

theorem identity_true (x y x' : a) :
    choose_bwd true x y x' = x' := rfl

theorem identity_false (x y y' : a) :
    choose_bwd false x y y' = y' := rfl
\end{lstlisting}

These four theorems constitute the \emph{forward-backward identity laws}:
applying the backward function with the forward result returns the input
unchanged. They are the functional analogue of the borrow-lifetime protocol
and suffice to reason about any function that selects a borrow based on
a pure condition.

\subsection{Comparison with Dijkstra Monads}

The LOOM~\cite{loom2026} and F$^\star$~\cite{fstar2016} approaches use Dijkstra
monads~\cite{dijkstraMonads2019}, which are indexed by pre- and post-conditions
in a way that generates verification conditions automatically. \RustM is
deliberately simpler: it is a plain monad without WP-index.

The choice reflects our use case. Dijkstra monads excel when the goal is to
automatically generate proof obligations from program code. Our goal is
different: we want to state explicit, human-readable specifications
(\ProgramSpec) and prove them by hand using Lean tactics. A plain monad
leaves the specification structure fully in the user's control and avoids
the complexity of a WP-monad universe hierarchy. Whether Dijkstra monads
would further automate the loop invariant proofs in \S\ref{sec:rust-proofs}
is an interesting open question we leave for future work.

%% ============================================================
%%  §4  The LLBC Deep Embedding
%% ============================================================

\section{The LLBC Deep Embedding}
\label{sec:llbc}

This section presents the deep embedding of LLBC into Lean~4. The embedding
consists of three parts: the AST types (\S\ref{sec:llbc-ast}), the interpreter
(\S\ref{sec:interpreter}), and the pipeline from Rust source to a verified
theorem (\S\ref{sec:pipeline}).

\subsection{The Verified Rust Fragment}
\label{sec:rust-fragment}

We verify a well-defined subset of safe Rust. The fragment is:

\begin{center}
\small
\begin{tabular}{@{}ll@{}}
\toprule
\textbf{Feature}             & \textbf{Status} \\
\midrule
Scalar integer types (I8/I16/I32/I64/U8/U16/U32/U64) & Supported \\
Boolean values               & Supported \\
Arithmetic with overflow checking & Supported \\
Comparisons, boolean logic   & Supported \\
Local variables (indexed)    & Supported \\
Assignments, sequencing      & Supported \\
\texttt{if}/\texttt{else}    & Supported \\
\texttt{loop}/\texttt{break}/\texttt{continue} & Supported \\
Function calls (cross-function via \texttt{FunEnv}) & Supported \\
\midrule
Struct fields, tuples        & Partial (projection in Charon; deref only in our AST) \\
\texttt{\&mut T} (mutable references) & Modelled via \texttt{FBPair} (\S\ref{sec:fbpair}) \\
\midrule
Raw pointers, \texttt{unsafe} blocks & \textbf{Not supported} \\
Closures, trait objects      & \textbf{Not supported} \\
Heap allocation (\texttt{Vec}, \texttt{Box}) & \textbf{Not supported} \\
Concurrency primitives       & \textbf{Not supported} \\
Lifetimes beyond \texttt{FBPair} & \textbf{Not supported} \\
\bottomrule
\end{tabular}
\end{center}

\noindent
This fragment matches what Charon can extract from safe Rust without
lifetimes or generic trait dispatch. Every function in the artifact
was checked to fall within this fragment before verification began.

\subsection{Trust Hierarchy}
\label{sec:tch}

The paper emphasises that the Lean kernel is the only trust anchor for
\emph{proof correctness}. However, the end-to-end pipeline introduces
additional trusted components for \emph{translation fidelity}:

\begin{center}
\small
\begin{tabular}{@{}lll@{}}
\toprule
\textbf{Component}    & \textbf{What it is trusted for} & \textbf{Verification status} \\
\midrule
Lean kernel           & All proof checks (258 theorems) & Formally verified \\
Mathlib               & Imported lemmas (\texttt{omega}, \texttt{Nat.fib}, etc.) & Community-reviewed \\
\midrule
Charon v0.1.197       & Rust MIR $\to$ LLBC JSON extraction & Unverified (external tool) \\
\texttt{charon2lean.py} & LLBC JSON $\to$ \texttt{LLBCFunDef} translation & Unverified (463 lines Python) \\
\bottomrule
\end{tabular}
\end{center}

\noindent
For the \textbf{hand-crafted} definitions in \texttt{ElabSpec.lean}
(F1--F16), the only components trusted are the Lean kernel and Mathlib:
the \texttt{LLBCFunDef} terms were written directly by the author and
inspected line by line against the Rust source. For \textbf{Charon-generated}
definitions (\texttt{CharonSpec.lean}, \texttt{NumIntegerSpec.lean}), the
trust chain extends to Charon and \texttt{charon2lean.py}: if either
mis-translates an instruction, the Lean kernel will verify a property about
an incorrect AST. This is analogous to the trusted code generator in Aeneas,
with the difference that our translator (Python, 463 lines) is much smaller
and auditable than a full compilation pipeline.

Mitigating this: (a)~every generated definition was manually reviewed against
the Charon JSON before proofs were written; (b)~ground-instance \rfl proofs
act as executable tests — if the AST is wrong, the kernel will not reduce to
the expected value.

\subsection{AST Types}
\label{sec:llbc-ast}

The LLBC type hierarchy in lean-pl-verify follows the structure of the Charon
intermediate representation~\cite{charon2024}. All types are defined as
inductive types in Lean~4.

\begin{lstlisting}[caption={LLBC type hierarchy.},label={lst:llbc-types}]
-- Scalar type tag (we support I32 and Bool in this fragment)
inductive LLBCTy where
  | I32 | Bool | Unit

-- A place is a local variable slot
inductive LLBCPlace where
  | var : Nat -> LLBCPlace

-- Operands are values that can be read from the current state
inductive LLBCOperand where
  | copy  : LLBCPlace -> LLBCOperand   -- borrow the value
  | move_ : LLBCPlace -> LLBCOperand   -- consume the value
  | const : LLBCLit   -> LLBCOperand   -- literal constant

-- Literals carry a type tag for scalar literals
inductive LLBCLit where
  | int  : Int -> LLBCTy -> LLBCLit
  | bool : Bool -> LLBCLit

-- Binary operators supported by the interpreter
inductive BinOp where
  | add | sub | mul | div | rem
  | eq | ne | lt | le | gt | ge
\end{lstlisting}

\texttt{LLBCPlace} is intentionally minimal: in the LLBC fragment we target
(safe Rust without raw pointers), every storage location is a numbered local
variable. The type supports extension to projections (\eg struct fields, tuple
components) but we leave that for future work.

\begin{lstlisting}[caption={LLBC rvalue and statement types.},label={lst:llbc-stmt}]
inductive LLBCRValue where
  | use    : LLBCOperand -> LLBCRValue
  | unOp   : UnOp -> LLBCOperand -> LLBCRValue
  | binOp  : BinOp -> LLBCOperand -> LLBCOperand -> LLBCRValue

inductive LLBCStmt where
  | assign   : LLBCPlace -> LLBCRValue -> LLBCStmt -> LLBCStmt
  | ite      : LLBCOperand -> LLBCStmt -> LLBCStmt -> LLBCStmt
  | loop     : LLBCStmt -> LLBCStmt
  | break_   : LLBCStmt
  | return_  : LLBCStmt
  | seq      : LLBCStmt -> LLBCStmt -> LLBCStmt
  | skip     : LLBCStmt

structure LLBCFunDef where
  name    : String
  params  : List LLBCTy
  locals  : List LLBCTy
  retTy   : LLBCTy
  body    : LLBCStmt

abbrev FunEnv := List LLBCFunDef
\end{lstlisting}

The \texttt{assign dst rvalue cont} constructor makes the continuation explicit:
the statement following an assignment is the third constructor field rather than
a separate sequencing node. This representation, adopted from CPS-style IRs,
slightly reduces the syntactic weight of typical function encodings; it means
that a function body like ``\texttt{x = a+b; y = x*2; return}'' is written
as a single nested \texttt{assign} rather than a chain of \texttt{seq} nodes.
Both styles are used in the artifact; \texttt{seq} is retained for the
post-loop continuation and explicit sequencing.

The \texttt{loop body} constructor takes a single body statement; the body
is expected to execute \texttt{break\_} to exit or fall through to restart.

\subsection{The Fuel-Parameterised Interpreter}
\label{sec:interpreter}

Lean~4's type theory requires all definitions to be terminating. LLBC loops
are potentially non-terminating, so the interpreter is parameterised by a
\emph{fuel} counter that decrements at each recursion step.

\begin{definition}[Fuel semantics]
\texttt{evalStmtFuel env fuel stmt locals} evaluates statement \texttt{stmt}
in function environment \texttt{env} with at most \texttt{fuel} recursive
steps, starting from local variable list \texttt{locals}. It returns
\texttt{.ok (sig, locals')} on success, where \texttt{sig} is a control-flow
signal (\texttt{.next}, \texttt{.break\_}, \texttt{.return\_}), or
\texttt{.error e} on panic.
\end{definition}

\begin{lstlisting}[caption={Core interpreter cases (simplified).},label={lst:interp}]
def evalStmtFuel (env : FunEnv) :
    Nat -> LLBCStmt -> Locals -> Except PanicReason (Signal × Locals)
  | _, .skip, s         => .ok (.next, s)
  | _, .break_, s       => .ok (.break_, s)
  | _, .return_, s      => .ok (.return_, s)

  -- Assignment: evaluate rvalue, update locals, run continuation
  | n+1, .assign dst rv k, s => do
      let v <- evalRValuePure rv s
      let s' := writePlacePure dst v s
      evalStmtFuel env n k s'

  -- Conditional: evaluate condition, dispatch to branch
  | n+1, .ite cond tb fb, s =>
      match evalOperandPure cond s with
      | some (.bool_ true)  => evalStmtFuel env n tb s
      | some (.bool_ false) => evalStmtFuel env n fb s
      | _ => .error .explicitPanic

  -- Loop: run body; .next restarts, .break_ exits
  | n+1, .loop body, s =>
      match evalStmtFuel env n body s with
      | .ok (.next,   s') => evalStmtFuel env n (.loop body) s'
      | .ok (.break_, s') => .ok (.next, s')
      | other             => other

  -- Explicit sequencing
  | n+1, .seq s1 s2, s =>
      match evalStmtFuel env n s1 s with
      | .ok (.next, s') => evalStmtFuel env n s2 s'
      | other           => other

  -- Out of fuel: treating this as a panic (not a logic error;
  -- correctness theorems ensure sufficient fuel is always supplied)
  | 0, _, _ => .error .explicitPanic
\end{lstlisting}

Two design decisions in this listing deserve explanation.

\textbf{Assignment continuation.} The \texttt{assign} case passes $n$ (not $n+1$)
to the continuation. This means each assignment step consumes one unit of fuel.
An intuitive way to think about this: fuel counts the number of evaluation
steps remaining, and each \texttt{assign} is one step.

\textbf{Loop fuel consumption.} In the \texttt{loop} case, the body is evaluated
with $n$ fuel (not $n+1$), and the recursive loop call also uses $n$ fuel.
This means the loop decrements by~1 per \emph{outer} loop invocation, not per
body step. The practical consequence is that a loop executing $k$ body iterations,
where each body needs $b$ fuel, requires $k \cdot b + k$ total fuel at the
outermost call. In the loop-correct theorems (\S\ref{sec:rust-proofs}), we
express this as $k + c \le \texttt{fuel}$ where $c$ accounts for the overhead
per iteration.

The helper functions \texttt{evalRValuePure} and \texttt{evalBinOpPure} are
pure (no monadic effects) and defined by case analysis; they return
\texttt{some v} on success and \texttt{none} on error. Binary operations check
overflow for \texttt{.add}, \texttt{.sub}, \texttt{.mul}:

\begin{lstlisting}[caption={evalBinOpPure for arithmetic.},label={lst:binop}]
def evalBinOpPure : BinOp -> Value -> Value -> Option Value
  | .add, .int a, .int b =>
      if minI32 <= a + b /\ a + b <= maxI32 then some (.int (a + b)) else none
  | .sub, .int a, .int b =>
      if minI32 <= a - b /\ a - b <= maxI32 then some (.int (a - b)) else none
  | .lt,  .int a, .int b => some (.bool_ (decide (a < b) = true))
  | .ge,  .int a, .int b => some (.bool_ (decide (a >= b) = true))
  -- ... further cases
\end{lstlisting}

The use of \texttt{decide (a < b) = true} rather than \texttt{a < b} is
deliberate: it produces a \texttt{Bool} via Lean's decidability mechanism,
which is kernel-evaluable on concrete integer values. The \texttt{Bool}-to-\texttt{Prop}
coercion is then handled by the \texttt{decide\_eq\_true} and
\texttt{decide\_eq\_false} lemmas in the conditional branch proofs.

\subsection{Function Evaluation}

\begin{lstlisting}[caption={Top-level function evaluation.},label={lst:evalfun}]
def buildLocals (f : LLBCFunDef) (args : List Value) : Locals :=
  -- allocate one slot for return value (index 0),
  -- then arg slots (index 1..params.length),
  -- then local variable slots (index params.length+1..end)
  let retSlot := [Value.unit]
  let argSlots := args.zipWith (fun _ v => v) f.params
  let localSlots := f.locals.map (fun _ => Value.unit)
  retSlot ++ argSlots ++ localSlots

def evalFunBody (env : FunEnv) (f : LLBCFunDef) (fuel : Nat)
    (locals : Locals) : Except PanicReason Value :=
  match evalStmtFuel env fuel f.body locals with
  | .ok (.return_, s) => (s[0]?).getD .unit |>.ok
  | .ok _             => .error .explicitPanic
  | .error e          => .error e

def evalFun (env : FunEnv) (f : LLBCFunDef) (fuel : Nat)
    (args : List Value) : RustM Unit Value :=
  fun _ =>
    match evalFunBody env f fuel (buildLocals f args) with
    | .ok v  => .ok (v, ())
    | .error e => .error e
\end{lstlisting}

The \texttt{at ()} syntax used in theorem statements (\eg
\texttt{evalFun [] f 30 args at ()}) is notation for applying the
\texttt{RustM Unit Value} computation to the trivial state \texttt{()}.
The state parameter is \texttt{Unit} because top-level function calls do not
share state with the caller; all state is local.

\subsection{Cross-Function Calls}
\label{sec:cross-calls}

The \texttt{FunEnv} list models a function environment that supports mutual
or cross-function calls. When a call instruction is evaluated, the interpreter
looks up the callee by name in \texttt{env} using \texttt{List.find?}:

\begin{lstlisting}[caption={Function call evaluation.},label={lst:funcall}]
  | n+1, .call fname args dest k, s =>
      match env.find? (fun f => f.name == fname) with
      | none   => .error .explicitPanic
      | some f =>
          let argVals := args.filterMap (fun op => evalOperandPure op s)
          match evalFunBody env f n argVals with
          | .ok v  => evalStmtFuel env n k (writePlacePure dest v s)
          | .error e => .error e
\end{lstlisting}

The cross-call uses $n$ fuel for the callee body, meaning the fuel is shared
between caller and callee steps. In the \texttt{square} $\to$ \texttt{mul}
example (\S\ref{sec:rust-proofs}), the concrete evaluation
\texttt{evalFun [mulFun] squareFun 20 [.int (-4)]} reduces by the
Lean kernel to \texttt{.ok (.int 16, ())} and the theorem is proved by \rfl.

\subsection{The Verification Pipeline}
\label{sec:pipeline}

The path from a Rust function to a Lean kernel-verified theorem proceeds
through five stages:

\begin{enumerate}[leftmargin=1.5em]
  \item \textbf{Write Rust source.} Functions are written in
        \texttt{examples/verified\_fns.rs}. Only safe Rust is used; the
        functions are simple enough that Charon can extract them without
        modifications.

  \item \textbf{Extract LLBC.} Run Charon on the Rust source to obtain the
        LLBC JSON.  The 463-line Python script \texttt{charon2lean.py}
        translates this JSON to \texttt{LLBCFunDef} Lean terms automatically
        (\S\ref{sec:eval-charon}).  Hand-crafted definitions in
        \texttt{ElabSpec.lean} (F1--F16) pre-date the pipeline and serve as
        the first verification layer; the pipeline is demonstrated on 18
        functions in \texttt{CharonSpec.lean}.

  \item \textbf{Encode as LLBCFunDef.} The LLBC AST is written as a Lean~4
        value of type \texttt{LLBCFunDef}. For example:

\begin{lstlisting}[caption={return42 function definition.},label={lst:r42}]
def return42Fun : LLBCFunDef := {
  name   := "return42"
  params := []
  locals := [.I32]
  retTy  := .I32
  body   := .assign (.var 0) (.use (.const (.int 42 .I32)))
              .return_
}
\end{lstlisting}

  \item \textbf{State the theorem.} The theorem asserts that \texttt{evalFun}
        on this definition satisfies the relevant \ProgramSpec. For
        \texttt{return42}:
\begin{lstlisting}
theorem elab_return42 :
    evalFun [] return42Fun 5 [] at () |=
      .pureOutput (. = .int 42) := by
  exact (.int 42, (), rfl, rfl)
\end{lstlisting}

  \item \textbf{Prove the theorem.} For pure functions, \rfl or
        $\langle\cdot, (), \texttt{rfl}, \texttt{rfl}\rangle$ closes the goal.
        For conditional functions, \texttt{simp only} with the relevant
        \texttt{decide\_eq\_true/false} lemma is used. For loops, the
        inductive invariant pattern described in \S\ref{sec:rust-proofs} is
        applied.
\end{enumerate}

This pipeline is entirely within Lean~4 and uses no external trusted components
beyond the Lean kernel and Mathlib. The 16 hand-crafted functions verified
(F1--F16) — \texttt{return42}, \texttt{id}, \texttt{neg}, \texttt{add},
\texttt{sub}, \texttt{max}, \texttt{abs}, \texttt{isZero}, \texttt{notGate},
\texttt{clamp}, \texttt{sumTo}, \texttt{fact}, \texttt{mul}, \texttt{min},
\texttt{square}, and \texttt{fib} — span the main categories of proof technique:
pure reduction (\rfl), overflow (\texttt{if\_pos}), conditional
(\texttt{decide\_eq\_true}), loop invariant (induction), and cross-function
(\rfl via kernel compute).

%% ============================================================
%%  §5  ProgramSpec: A Shared Specification Language
%% ============================================================

\section{ProgramSpec: A Shared Specification Language}
\label{sec:programspec}

A central claim of this work is that a single specification language can
describe properties of programs from different source languages. This section
defines that language — \ProgramSpec — and demonstrates that its 17 proved
combinators are sufficient to state and verify all correctness properties in
the library.

\subsection{Motivation}

Consider two implementations of \texttt{max(a, b)}:
\begin{enumerate}[leftmargin=1.5em]
  \item A Rust function compiled to an \texttt{LLBCFunDef} and run via
        \texttt{evalFun}.
  \item A TypeScript function compiled to a \texttt{TSFunDef} and run via
        \texttt{evalTSFun}.
\end{enumerate}

Both should satisfy the same specification: ``if $a \ge b$, the result is
$a$; otherwise the result is $b$.'' Writing this specification once and
applying it to both implementations requires a specification language that
is \emph{parametric in the execution mechanism}. \ProgramSpec achieves this
by specifying properties of \RustM computations, which both evaluators return.

\subsection{The ProgramSpec Inductive}

\begin{lstlisting}[caption={ProgramSpec inductive type.},label={lst:progspec}]
inductive ProgramSpec (s a : Type) where
  | pureOutput : (a -> Prop) -> ProgramSpec s a
  | nocrash    : ProgramSpec s a
  | both       : ProgramSpec s a -> ProgramSpec s a -> ProgramSpec s a
  | withFuel   : Nat -> ProgramSpec s a -> ProgramSpec s a
  | precond    : Prop -> ProgramSpec s a -> ProgramSpec s a
\end{lstlisting}

The constructors have the following meanings:

\begin{itemize}[leftmargin=1.5em]
  \item \texttt{pureOutput P}: the computation returns some value $v$ such
        that $P\,v$ holds and does not modify the state.
  \item \texttt{nocrash}: the computation does not return \texttt{.error}.
  \item \texttt{both P Q}: both $P$ and $Q$ are satisfied simultaneously.
  \item \texttt{withFuel n P}: $P$ is satisfied when the computation is
        given at least $n$ units of fuel. (This constructor is used for
        fuel-parametric specifications and is not used in the main theorem
        statements.)
  \item \texttt{precond h P}: if the proposition $h$ holds, then $P$ is
        satisfied. This models preconditioned specifications.
\end{itemize}

\subsection{The Satisfaction Relation}

The satisfaction relation $m \satisfies P$ is defined by structural recursion
on $P$:

\begin{lstlisting}[caption={Satisfaction relation.},label={lst:sat}]
def satisfies (m : RustM s a) (init : s) :
    ProgramSpec s a -> Prop
  | .pureOutput P => exists  v, m init = .ok (v, init) /\ P v
  | .nocrash      => forall  e, m init /= .error e
  | .both P Q     => satisfies m init P /\ satisfies m init Q
  | .withFuel n P => satisfies m init P  -- simplified form
  | .precond h P  => h -> satisfies m init P

-- Notation: m at s |= P  means satisfies m s P
notation:50 m " at " s " |= " P => satisfies m s P
\end{lstlisting}

The \texttt{pureOutput} case deserves attention. The condition
\texttt{m init = .ok (v, init)} (not \texttt{.ok (v, s')} for some
\emph{different} $s'$) requires the state to be unchanged. This is always
satisfied by top-level function calls (where the state is \texttt{Unit}),
but it would fail for computations that write to a shared heap.
Since we do not model a shared heap, this simplification is sound.

\subsection{The Combinator Library}

The \texttt{Spec/Satisfies.lean} module proves 17 combinators. We list the
most frequently used ones:

\begin{lstlisting}[caption={Key ProgramSpec combinators (selection).},label={lst:combinators}]
-- Pure return: pure a satisfies pureOutput (. = a)
theorem sat_pure (a : a) (s : s) :
    pure a at s |= .pureOutput (. = a) := by
  simp [satisfies, pure_ok]

-- Sequential composition: if f returns a, and g a satisfies P, then f >>= g satisfies P
theorem sat_bind_pureOutput {m : RustM s a} {f : a -> RustM s b}
    (hm : m at s |= .pureOutput (. = a))
    (hf : f a at s |= P) :
    (m >>= f) at s |= P := by
  obtain (v, hv, rfl) := hm
  simp [bind_ok hm ...]

-- Conjunction introduction
theorem sat_both_intro
    (hP : m at s |= P) (hQ : m at s |= Q) :
    m at s |= .both P Q := (hP, hQ)

-- Monotonicity: if (. = v) implies P, then pureOutput (. = v) implies pureOutput P
theorem sat_pureOutput_mono {m : RustM s a}
    (h  : m at s |= .pureOutput (. = v))
    (hP : P v) :
    m at s |= .pureOutput P := by
  obtain (w, hw, rfl) := h; exact (w, hw, hP)
\end{lstlisting}

\texttt{sat\_pureOutput\_mono} is particularly important for the TypeScript
theorems (\S\ref{sec:typescript}): when the base-case execution equation is
blocked by the \texttt{Int.decLe @[extern]} issue, we can still derive
higher-level specifications from lower-level ones via this monotonicity lemma,
avoiding the need for \sorry in the derived theorems.

\begin{table}[t]
\centering
\caption{All 17 ProgramSpec combinators.}
\label{tab:combinators}
\small
\begin{tabular}{@{}ll@{}}
\toprule
\textbf{Combinator} & \textbf{Informal statement} \\
\midrule
\texttt{sat\_pure}              & \texttt{pure a} satisfies \texttt{pureOutput (. = a)} \\
\texttt{sat\_bind\_pureOutput}  & Compose: if $m$ returns $a$, then $m >>= f$ satisfies $f$'s spec \\
\texttt{sat\_both\_intro}       & Introduce conjunction \\
\texttt{sat\_both\_elim\_left}  & Eliminate left conjunct \\
\texttt{sat\_both\_elim\_right} & Eliminate right conjunct \\
\texttt{sat\_pureOutput\_mono}  & Weaken: \texttt{(.=v)} $\Rightarrow$ $P\,v$ yields $P$ \\
\texttt{sat\_nocrash\_of\_ok}   & If $m$ returns \texttt{ok}, it does not crash \\
\texttt{sat\_precond\_intro}    & Prove preconditioned spec by providing the precondition \\
\texttt{sat\_precond\_elim}     & Apply preconditioned spec \\
\texttt{sat\_rpanic\_nocrash}   & $\neg$: \texttt{rpanic} violates \texttt{nocrash} \\
\texttt{sat\_pureOutput\_eq}    & \texttt{pureOutput (. = v)} is extensionally equal \\
\texttt{sat\_map\_pureOutput}   & Functor: map preserves \texttt{pureOutput} \\
\texttt{sat\_bind\_nocrash}     & Bind: if $m$ and $f$ both don't crash, neither does $m >>= f$ \\
\texttt{sat\_precond\_True}     & If precondition is \texttt{True}, it can be discharged \\
\texttt{sat\_both\_comm}        & Commutativity of \texttt{both} \\
\texttt{sat\_withFuel\_intro}   & Introduce fuel-parametric spec \\
\texttt{sat\_withFuel\_elim}    & Eliminate fuel-parametric spec \\
\bottomrule
\end{tabular}
\end{table}

\subsection{RustM-Level Examples}

Before connecting \ProgramSpec to the LLBC embedding, it is useful to see it
applied to hand-written monadic terms. The \texttt{Spec/Examples.lean} module
contains eight such examples (E1--E8) to validate the combinator library in
isolation from the interpreter.

\begin{lstlisting}[caption={E1: rustAdd satisfies its specification.},label={lst:e1}]
-- Direct monadic definition of addition with overflow check
def rustAdd (x y : Int) : RustM Unit Int :=
  if minI32 <= x + y /\ x + y <= maxI32 then pure (x + y)
  else rpanic .arithmeticOverflow

theorem e1_add_correct (x y : Int)
    (h : minI32 <= x + y /\ x + y <= maxI32) :
    rustAdd x y at () |= .pureOutput (. = x + y) := by
  simp [rustAdd, satisfies, if_pos h]
\end{lstlisting}

\begin{lstlisting}[caption={E5: rustMax with two-case specification.},label={lst:e5}]
def rustMax (a b : Int) : RustM Unit Int :=
  pure (if a >= b then a else b)

theorem e5_max_ge (a b : Int) (h : a >= b) :
    rustMax a b at () |= .pureOutput (. = a) := by
  simp [rustMax, satisfies, if_pos h]

theorem e5_max_lt (a b : Int) (h : ¬(a >= b)) :
    rustMax a b at () |= .pureOutput (. = b) := by
  simp [rustMax, satisfies, if_neg h]

theorem e5_max_both (a b : Int) :
    rustMax a b at () |=
      .both (.pureOutput (. = if a >= b then a else b))
            .nocrash := by
  refine sat_both_intro ?_ ?_
  . simp [rustMax, satisfies]
  . simp [rustMax, satisfies]
\end{lstlisting}

These examples exercise the full combinator library and establish confidence
that the specification language is adequate before the LLBC interpreter is
introduced.

\subsection{Language Independence}

The \ProgramSpec type contains no reference to \texttt{LLBCFunDef},
\texttt{LLBCStmt}, \texttt{evalStmtFuel}, \texttt{TSFunDef}, or any other
language-specific construct. It is parametric in $\sigma$ (the state type)
and $\alpha$ (the value type). The satisfaction relation applies to any
\texttt{RustM}~$\sigma$~$\alpha$ computation.

This parametricity is the mechanism by which cross-language unification
theorems become possible. When the Rust evaluator \texttt{evalFun} and the
TypeScript evaluator \texttt{evalTSFun} both return computations of type
\texttt{RustM Unit Value}, the same \ProgramSpec specification
\texttt{.pureOutput (. = .int (if a $\ge$ b then a else b))} can be stated
and proved for both, and the two proofs can be composed into a unification
theorem without any additional logical infrastructure.

%% ============================================================
%%  §6  Verification of Rust Programs via LLBC
%% ============================================================

\section{Verification of Rust Programs via LLBC}
\label{sec:rust-proofs}

The \texttt{Translation/ElabSpec.lean} module contains 48 theorems covering
16 Rust functions. This section describes the four proof strategies used,
then gives a detailed walkthrough of the loop invariant methodology using
\texttt{sumTo} as the primary example.

\subsection{Theorem Shape}

All LLBC theorems have the same outer form:

\begin{lstlisting}[caption={Standard LLBC theorem shape.},label={lst:thmshape}]
theorem elab_f_correct (args...) (precond...) :
    evalFun env fFun fuel args at () |=
      .pureOutput (. = expected_value) := ...
\end{lstlisting}

The \texttt{at ()} applies the \texttt{RustM Unit Value} computation to the
trivial unit state. Unfolding \texttt{satisfies}, this expands to:
$\exists v,\; \texttt{evalFun env fFun fuel args () = .ok (v, ())} \wedge v = \text{expected}$.

In practice, the proof witnesses $v$ directly and proves the two conjuncts
separately: the first by kernel computation (often \rfl), the second by
definition unfolding.

\subsection{Pure Functions: Proof by \rfl}

For functions with no conditional branches and no arithmetic overflow, the
entire evaluation reduces by kernel computation:

\begin{lstlisting}[caption={F1: return42 is proved by rfl.},label={lst:r42proof}]
theorem elab_return42 :
    evalFun [] return42Fun 5 [] at () |=
      .pureOutput (. = .int 42) :=
  (.int 42, (), rfl, rfl)
\end{lstlisting}

The \rfl in the first conjunct position asks the kernel to verify that
\texttt{evalFun [] return42Fun 5 [] () = .ok (.int 42, ())}.
The kernel unfolds \texttt{evalFun}, \texttt{evalFunBody}, \texttt{evalStmtFuel},
\texttt{evalRValuePure}, \texttt{buildLocals}, and the \texttt{return42Fun}
definition, reducing the entire expression to \texttt{.ok (.int 42, ())} in
a few hundred reduction steps.

This same strategy works for \texttt{id}, \texttt{neg}, and any function
whose body is a finite sequence of unconditional assignments. The key
requirement is that no \texttt{@[extern]} annotation blocks the reduction.

\subsection{Arithmetic with Overflow: \texttt{if\_pos} Pattern}

For functions that perform checked arithmetic, the interpreter's overflow guard
introduces a conditional that cannot be resolved by kernel reduction alone:

\begin{lstlisting}[caption={F4: add with overflow precondition.},label={lst:addproof}]
theorem elab_add (x y : Int)
    (h : minI32 <= x + y /\ x + y <= maxI32) :
    evalFun [] addFun (10) [.int x, .int y] at () |=
      .pureOutput (. = .int (x + y)) := by
  refine (.int (x + y), (), ?_, rfl)
  simp only [evalFun, evalFunBody, evalStmtFuel, buildLocals,
    evalRValuePure, evalBinOpPure, evalOperandPure,
    evalPlacePure, writePlacePure, List.set, List.replicate]
  rw [if_pos h]
  rfl
\end{lstlisting}

The \texttt{simp only} reduces all of the interpreter's helper functions
except the overflow guard. The \texttt{rw [if\_pos h]} resolves the guard
using the hypothesis $h$, after which \rfl closes the goal.

This pattern — \texttt{simp only [interpreter lemmas]; rw [if\_pos h]; rfl}
— applies uniformly to \texttt{sub} and \texttt{mul} as well. The only
variation is the precondition $h$.

\subsection{Conditionals: \texttt{decide\_eq\_true} Pattern}

Comparison operations produce \texttt{.bool\_ (decide (a $\ge$ b) = true)},
which requires a different reduction approach:

\begin{lstlisting}[caption={F6: max function, ge branch.},label={lst:maxproof}]
theorem elab_max_ge (a b : Int) (h : a >= b) :
    evalFun [] maxFun 15 [.int a, .int b] at () |=
      .pureOutput (. = .int a) := by
  refine (.int a, (), ?_, rfl)
  have hd : decide (a >= b) = true := decide_eq_true h
  simp only [evalFun, ..., maxFun, hd, if_true]

theorem elab_max_lt (a b : Int) (h : ¬(a >= b)) :
    evalFun [] maxFun 15 [.int a, .int b] at () |=
      .pureOutput (. = .int b) := by
  refine (.int b, (), ?_, rfl)
  have hd : decide (a >= b) = false := decide_eq_false h
  simp only [evalFun, ..., maxFun, hd, if_false]

-- Derived: both-case theorem using sat_pureOutput_mono
theorem elab_max_either (a b : Int) :
    evalFun [] maxFun 15 [.int a, .int b] at () |=
      .pureOutput (. = .int (if a >= b then a else b)) := by
  by_cases h : a >= b
  . exact sat_pureOutput_mono (elab_max_ge a b h) (by simp [h])
  . exact sat_pureOutput_mono (elab_max_lt a b h) (by simp [h])
\end{lstlisting}

The two-theorem pattern (one per branch) plus a derived ``either'' theorem
is used for all conditional functions: \texttt{max}, \texttt{min}, \texttt{abs},
\texttt{isZero}, \texttt{notGate}. The derived theorem is the most useful for
consumers of the API; the branch-specific theorems are stepping stones.

\subsection{Loop Invariant Methodology: The \texttt{sumTo} Walkthrough}

The \texttt{sumTo(n)} function computes $\sum_{i=0}^{n-1} i = n(n-1)/2$
(the Gauss formula) using a loop. This is the primary running example for
the loop invariant methodology. Four components are needed:

\begin{enumerate}[leftmargin=2em,label=\textbf{\arabic*.}]
  \item A \emph{specification function} $\texttt{sumRange}(i, n)$
        computing the partial sum $\sum_{k=i}^{n-1} k$.
  \item An \emph{overflow predicate} \SumToOv asserting that, at every
        remaining iteration, the accumulator stays within I32 bounds.
  \item A \emph{loop-correct theorem} proved by induction on the remaining
        iteration count.
  \item A \emph{safety-implies-overflow bridge lemma} establishing that
        safe inputs (here, $0 \le n \le 46\,340$) imply the overflow predicate.
\end{enumerate}

\subsubsection{Specification Function}

\begin{lstlisting}[caption={sumRange specification.},label={lst:sumrange}]
def sumRange : Int -> Int -> Int
  | i, n => if i >= n then 0
             else i + sumRange (i + 1) n
termination_by n - i

theorem sumRange_gaussSum (n : Int) (hn : 0 <= n) :
    sumRange 0 n = n * (n - 1) / 2 := by
  induction n with
  | ofNat k => induction k with ...
  | negSucc k => omega
\end{lstlisting}

\subsubsection{Overflow Predicate}

\begin{lstlisting}[caption={SumToOv: overflow condition for sumTo.},label={lst:sumtoov}]
def SumToOv (n acc i : Int) : Prop :=
  forall  j : Int, i <= j -> j < n ->
    (minI32 <= acc + sumRange i j /\
     acc + sumRange i j <= maxI32) /\
    (minI32 <= j + 1 /\ j + 1 <= maxI32)
\end{lstlisting}

This predicate is parametric in the current loop counter $i$ and the
accumulator $\texttt{acc}$. At each iteration $j \in [i, n)$, it asserts that
the accumulator at step $j$ (which is $\texttt{acc} + \texttt{sumRange}(i, j)$)
stays in range, and the loop counter does not overflow.

The critical step lemma says the predicate is preserved by one iteration:

\begin{lstlisting}[caption={SumToOv_step: invariant preservation.},label={lst:sumtoovstep}]
theorem SumToOv_step {n acc i : Int}
    (hlt : i < n)
    (hov : SumToOv n acc i) :
    SumToOv n (acc + i) (i + 1) := by
  intro j hj hjn
  have := hov j (by omega) hjn
  simpa [sumRange, hlt] using this
\end{lstlisting}

\subsubsection{Loop-Correct Theorem}

\begin{lstlisting}[caption={sumTo_loop_correct: inductive invariant.},label={lst:sumtolc}]
theorem sumTo_loop_correct
    (k : Nat) (env : FunEnv)
    (n acc i : Int) (c3 c4 : Value)
    (heq  : (n - i).toNat = k)
    (hi   : 0 <= i) (hin : i <= n)
    (hinv : acc = sumRange 0 i)   -- invariant: acc accumulates sum so far
    (hov  : SumToOv n acc i)
    (fuel : Nat) (hfuel : k + 8 <= fuel) :
    exists  c3' c4',
    evalStmtFuel env fuel (.loop sumToBody)
      [.unit, .int n, .int acc, .int i, c3, c4] =
    .ok (.next, [.unit, .int n,
                 .int (sumRange 0 n), .int n, c3', c4']) := by
  induction k generalizing n acc i c3 c4 fuel with
  | zero =>
    -- i = n: loop exits, acc = sumRange 0 n
    have hieqn : i = n := le_antisymm hin (by omega)
    ...
  | succ k' ih =>
    -- i < n: one iteration advances i and adds i to acc
    have hlt : i < n := by omega
    ...
    exact ih n (acc + i) (i + 1) ... (SumToOv_step hlt hov) ...
\end{lstlisting}

The induction variable is $k = (n - i).\texttt{toNat}$, the number of
remaining iterations. The zero case ($k = 0$, \ie $i = n$) uses the
\texttt{sumTo\_body\_false} step lemma to show the loop exits with the
accumulator in its final state. The successor case ($k = k' + 1$) applies
\texttt{sumTo\_body\_true} to advance one iteration, then calls the
induction hypothesis with the updated accumulator and counter.

\subsubsection{The \texttt{hpre} Technique}

A subtle proof engineering challenge arises in the top-level theorem. The
function body includes three prefix assignments (initialising \texttt{acc} and
\texttt{i}) before the loop, then a postfix assignment reading the return value.
If we unfold \texttt{evalFun} naively, the goal becomes:

\begin{center}
\small\texttt{evalStmtFuel env fuel (assign ... (assign ... (assign ... (loop sumToBody; assign ... return\_))))}
\end{center}

If the \texttt{simp} tactic attempts to simplify this as a whole, it will
unfold the loop construct too eagerly, losing the structured induction.

The solution is the \texttt{hpre} lemma pattern: we first prove a
``prefix reduction'' fact by \rfl, which shows that the full body reduces to
the loop with known initial state, \emph{without} unfolding the loop:

\begin{lstlisting}[caption={hpre technique for prefix reduction.},label={lst:hpre}]
have hpre : evalStmtFuel env (n.toNat + fuel_overhead) sumToFun.body
    (buildLocals sumToFun [.int n]) =
    (match evalStmtFuel env (n.toNat + k_overhead) (.loop sumToBody)
         [.unit, .int n, .int 0, .int 0, .unit, .unit] with
     | .ok (.next, s') =>
         evalStmtFuel env ... (.assign (.var 0) (.use (.copy (.var 2))) .return_) s'
     | other => other) := rfl
\end{lstlisting}

The \rfl succeeds because the Lean kernel can evaluate the prefix assignments
and the \texttt{seq} combinator, stopping at the \texttt{loop} constructor.
Once \texttt{hpre} is established, \texttt{simp only [evalFun, hpre, hloop]}
closes the goal cleanly.

\subsection{Factorial: Completing the Loop Invariant}

The \texttt{fact} function's loop invariant follows the same pattern as
\texttt{sumTo}, with \texttt{FactOv} replacing \SumToOv and
\texttt{prodRange} replacing \texttt{sumRange}. All components are
fully proved (0 \sorry): \texttt{fact\_body\_true}, \texttt{fact\_body\_false},
\texttt{FactOv\_step}, \texttt{fact\_loop\_correct}, and the two auxiliary
lemmas below.

Two auxiliary lemmas required particular proof engineering:

\begin{itemize}[leftmargin=1.5em]
  \item \texttt{prodRange\_factorial}: $\prod_{k=1}^{n} k = n!$, proved by
        induction on $n$ using a right-peel lemma
        (\texttt{prodRange\_succ}) that decomposes $\prod_{k=1}^{n+1} k$
        into $\prod_{k=1}^{n} k \cdot (n+1)$, closed by \texttt{omega}.
  \item \texttt{factSafe\_implies\_ov}: for $0 \le j \le 12$,
        $j! \le 12! < 2^{31} - 1$. Proved by \texttt{interval\_cases j}
        (a Mathlib tactic that exhaustively cases over the 13 values)
        followed by \texttt{norm\_num} at each leaf.
\end{itemize}

Both lemmas are verified by the Lean kernel with \textbf{0 \sorry}.
The \texttt{FactInvariant.lean} file contributes 9 theorems to the total.

\subsection{Ground Instances}

For fixed inputs, every loop function reduces by kernel computation:

\begin{lstlisting}[caption={Ground instances proved by rfl.},label={lst:ground}]
-- sumTo(5) = 10 (Gauss formula: 5*4/2 = 10)
theorem elab_sumTo_five :
    evalFun [] sumToFun 20 [.int 5] at () |=
      .pureOutput (. = .int 10) :=
  (.int 10, (), rfl, rfl)

-- fact(5) = 120
theorem elab_fact_five :
    evalFun [] factFun 60 [.int 5] at () |=
      .pureOutput (. = .int 120) :=
  (.int 120, (), rfl, rfl)

-- fib(10) = 55
theorem elab_fib_ten :
    evalFun [] fibFun 25 [.int 10] at () |=
      .pureOutput (. = .int 55) :=
  (.int 55, (), rfl, rfl)
\end{lstlisting}

These are not trivial: the Lean kernel is literally executing the loop —
all 5 iterations of \texttt{sumTo(5)}, all 5 iterations of \texttt{fact(5)},
all 10 iterations of \texttt{fib(10)} — and confirming the result. The fact
that \rfl succeeds is direct evidence that the interpreter is correctly
encoding the loop semantics.

%% ============================================================
%%  §6b  Relational Semantics and Interpreter Adequacy
%% ============================================================

\section{Relational Semantics and Interpreter Adequacy}
\label{sec:semantics}

The \texttt{Translation/ElabSpec.lean} theorems (\S\ref{sec:rust-proofs})
prove that the \emph{fuel-based interpreter} produces correct outputs.
A natural question is whether the interpreter faithfully captures the
\emph{intended meaning} of LLBC programs — i.e., whether it is adequate
with respect to a language-independent semantic ground truth.
This section answers that question via a relational big-step semantics
(\texttt{Translation/Semantics.lean}) and a soundness proof
(\texttt{Translation/Adequacy.lean}).

\subsection{Relational Big-Step Semantics}

We define a standard inductive big-step operational semantics over the
same runtime types (\texttt{Value}, \texttt{Locals}, \texttt{Signal},
\texttt{FunEnv}) as the interpreter — sharing types avoids any
translation layer in the adequacy proof.

The semantics consists of six mutually inductive relations:

\begin{center}
\small
\begin{tabular}{@{}lll@{}}
\toprule
\textbf{Relation} & \textbf{Judgment} & \textbf{Meaning} \\
\midrule
\texttt{EvalPlace}   & $s \vdash p \Downarrow v$           & Place $p$ reads value $v$ from locals $s$ \\
\texttt{EvalOperand} & $s \vdash o \Downarrow v$           & Operand $o$ evaluates to $v$ \\
\texttt{EvalRValue}  & $s \vdash r \Downarrow v$           & R-value $r$ evaluates to $v$ \\
\texttt{EvalWrite}   & $s \xrightarrow{p := v} s'$         & Writing $v$ to place $p$ yields $s'$ \\
\texttt{EvalStmt}    & $\Gamma \vdash s, \sigma \Downarrow \sigma', s'$ & Statement executes with signal $\sigma$ \\
\texttt{EvalFun}     & $\Gamma \vdash f(\overline{v}) \Downarrow v'$ & Function call terminates with $v'$ \\
\bottomrule
\end{tabular}
\end{center}

\noindent
\textbf{Design choices.}
Loops are handled by three rules — \texttt{loop\_cont} (body produces
\texttt{next}; continue looping), \texttt{loop\_break} (body produces
\texttt{break\_}; exit), and \texttt{loop\_return} (body produces
\texttt{return\_}; propagate) — so the relation stays strictly inductive
(no coinduction).
Panics are modelled by \emph{absence} of a derivation: there is no
$\texttt{EvalStmt}$ rule for division by zero or integer overflow.
This matches the Lean convention that partial functions are undefined
on bad inputs rather than mapped to an error value.

\begin{lstlisting}[caption={Selected \texttt{EvalStmt} rules (Lean notation).},
                   label={lst:evalstmt}]
-- Sequential composition: s1 produces next, then s2 runs
| seq_next {s1 s2 sig s s_mid s'} :
    EvalStmt env s1 s .next s_mid →
    EvalStmt env s2 s_mid sig s' →
    EvalStmt env (.seq s1 s2) s sig s'

-- Loop continues: body produces .next
| loop_cont {body s s' s''} :
    EvalStmt env body s .next s' →
    EvalStmt env (.loop body) s' sig s'' →
    EvalStmt env (.loop body) s sig s''

-- Loop breaks: body produces .break_
| loop_break {body s s'} :
    EvalStmt env body s .break_ s' →
    EvalStmt env (.loop body) s .next s'
\end{lstlisting}

\subsection{Soundness}

The main adequacy theorem asserts that every output produced by the
fuel-based interpreter is derivable in the relational semantics:

\begin{lstlisting}[caption={Soundness of the LLBC interpreter.},
                   label={lst:soundness}]
theorem evalStmtFuel_sound
    {env : FunEnv} {n : Nat} {stmt : LLBCStmt}
    {s : Locals} {sig : Signal} {s' : Locals}
    (h : evalStmtFuel env n stmt s = Except.ok (sig, s')) :
    EvalStmt env stmt s sig s' := ...
\end{lstlisting}

\noindent
The proof is by structural induction on $n$ (the fuel counter), with a
case split on the statement constructor.
Each case reduces the hypothesis \texttt{h} to a concrete
\texttt{Except.ok} equation via \texttt{simp} and \texttt{split},
then constructs the corresponding relational derivation.
The proof comprises 14 theorems (0 sorry), covering all LLBC statement forms:
\texttt{assign}, \texttt{return\_}, \texttt{skip}, \texttt{break\_},
\texttt{continue\_}, \texttt{seq} (both next and abort variants),
\texttt{ite}, \texttt{loop} (all three exit modes), and \texttt{call}.

\paragraph{Key proof patterns.}
Two patterns appear throughout the adequacy proof:

\begin{itemize}[leftmargin=2em]
  \item \textbf{\texttt{split at h}} for nested \texttt{match} reductions.
        Using \texttt{rcases} + \texttt{simp} instead leads to lambda
        alpha-renaming issues that break dependent pattern matching; \texttt{split}
        avoids this by keeping the match scrutinee in the context.

  \item \textbf{\texttt{change} for do-notation} reductions.
        Lean's kernel reduces \texttt{do}-notation to nested \texttt{match} on
        \texttt{Except.ok} / \texttt{Except.error}; the \texttt{change}
        tactic asserts the reduced form explicitly, avoiding opaque unfolding
        steps that would otherwise require \texttt{norm\_num} or \texttt{ring}.
\end{itemize}

\subsection{Completeness}

The converse direction — every relational derivation is witnessed by
some fuel bound — is proved as \texttt{evalStmtFuel\_complete}:

\begin{lstlisting}[caption={Completeness of the LLBC interpreter.},
                   label={lst:complete}]
theorem evalStmtFuel_complete
    {env : FunEnv} {stmt : LLBCStmt}
    {s : Locals} {sig : Signal} {s' : Locals}
    (h : EvalStmt env stmt s sig s') :
    ∃ n, evalStmtFuel env n stmt s = Except.ok (sig, s') := ...
\end{lstlisting}

\noindent
The proof proceeds by structural induction on the \texttt{EvalStmt} derivation.
Two auxiliary lemmas are key:

\begin{itemize}[leftmargin=2em]
  \item \textbf{\texttt{evalStmtFuel\_mono\_one}}: if
        \texttt{evalStmtFuel env n stmt s = .ok r}, then
        \texttt{evalStmtFuel env (n+1) stmt s = .ok r}.
        Proved by induction on $n$, case analysis on \texttt{stmt}.

  \item \textbf{\texttt{evalStmtFuel\_mono}}: monotonicity for all $k$,
        derived from \texttt{evalStmtFuel\_mono\_one} by induction on $k$.
\end{itemize}

\noindent
For the \texttt{seq\_next} case, the existential witnesses are
$\max(n_1, n_2) + 1$ where $n_1, n_2$ come from the two inductive
hypotheses; \texttt{evalStmtFuel\_mono} is used to lift each sub-proof
to the common bound.
The \texttt{loop\_cont} case requires composing the monotonicity lemma for
the body and the recursive loop call; this is the most technically intricate
case but follows the same pattern.

Together, \texttt{evalStmtFuel\_sound} and \texttt{evalStmtFuel\_complete}
establish the \emph{full adequacy} of the interpreter: the fuel-based
operational semantics and the inductive big-step semantics are extensionally
equivalent. The \texttt{Translation/Adequacy.lean} file contains
\textbf{19 theorems (0 \sorry)}.

\subsection{Relationship to the Verification Pipeline}

The relational semantics serves as the \emph{semantic ground truth} against
which the interpreter is calibrated:

\[
  \underbrace{\texttt{evalStmtFuel}}_{\text{interpreter}}
  \;\overset{\text{sound}}{\longrightarrow}\;
  \underbrace{\texttt{EvalStmt}}_{\text{relational semantics}}
\]

The \texttt{ElabSpec} and \texttt{CharonSpec} theorems operate at the
interpreter level. Soundness guarantees that any property proved there
also holds in the relational model. This closes the semantic gap between
the constructive evaluation approach (easy to compute with, easy to prove
concrete instances by \texttt{rfl}) and the declarative inductive semantics
(easy to reason about symbolically, standard in PL theory).

%% ============================================================
%%  §7  Proof Automation
%% ============================================================

\section{Proof Automation}
\label{sec:automation}

The 48 LLBC theorems in \texttt{ElabSpec.lean} share a heavily redundant
proof structure. A typical pure-function proof rewrites through the same
list of 12--15 simp lemmas; a typical conditional-function proof repeats
the \texttt{simp}/\texttt{rw [if\_pos]}/\rfl sequence. Writing each proof
separately is verbose, and the repeated simp sets are brittle: adding a
new interpreter helper function requires updating every proof.

This section describes the four tactic macros that address this problem,
the non-obvious design issues encountered, and the compression ratio achieved.

\subsection{The Fixed Simp Set}

All LLBC evaluation proofs share a core set of simp lemmas:

\begin{lstlisting}[caption={The shared LLBC simp set.},label={lst:simpset}]
-- These 12 lemmas unfold the interpreter on any closed LLBC term
private def llbcSimp := [
  evalFun, evalFunBody, evalStmtFuel, buildLocals,
  evalRValuePure, evalBinOpPure, evalOperandPure,
  evalPlacePure, writePlacePure,
  List.set, List.replicate, List.foldl
]
\end{lstlisting}

In a standalone proof, this list appears explicitly in every \texttt{simp only
[\ldots]} call, typically spanning 2--3 lines. With 48 theorems, that is
roughly 100--150 lines of simp-set boilerplate.

\subsection{The Four Macros}

\subsubsection{llbc\_verify: Pure Functions}

\begin{lstlisting}[caption={llbc\_verify macro.},label={lst:lv}]
macro "llbc_verify" f:term : tactic =>
  `(tactic| exact (_, (), rfl, rfl) <|>
    (simp only [evalFun, evalFunBody, evalStmtFuel, buildLocals,
      evalRValuePure, evalBinOpPure, evalOperandPure,
      evalPlacePure, writePlacePure, List.set,
      List.replicate, List.foldl, $f:term]))
\end{lstlisting}

Used for pure functions (\texttt{return42}, \texttt{id}, \texttt{neg}):

\begin{lstlisting}
theorem elab_return42 : evalFun [] return42Fun 5 [] at () |=
    .pureOutput (. = .int 42) := by
  llbc_verify return42Fun
\end{lstlisting}

\subsubsection{llbc\_verify\_prop: Arithmetic with Overflow}

\begin{lstlisting}[caption={llbc\_verify\_prop macro.},label={lst:lvp}]
syntax "llbc_verify_prop" ident term : tactic
macro_rules
  | `(tactic| llbc_verify_prop $f $h) =>
    `(tactic|
      refine (_, (), ?_, rfl)
      simp only [evalFun, evalFunBody, evalStmtFuel, buildLocals,
        evalRValuePure, evalBinOpPure, evalOperandPure,
        evalPlacePure, writePlacePure, List.set,
        List.replicate, List.foldl, $f:ident]
      rw [if_pos $h:term]
      rfl)
\end{lstlisting}

Used for arithmetic functions (\texttt{add}, \texttt{sub}, \texttt{mul}):

\begin{lstlisting}
theorem elab_add (x y : Int) (h : minI32 <= x + y /\ x + y <= maxI32) :
    evalFun [] addFun 10 [.int x, .int y] at () |=
      .pureOutput (. = .int (x + y)) := by
  llbc_verify_prop addFun h
\end{lstlisting}

\subsubsection{llbc\_verify\_cond: Conditional Functions}

\begin{lstlisting}[caption={llbc\_verify\_cond macro.},label={lst:lvc}]
syntax "llbc_verify_cond" ident term : tactic
macro_rules
  | `(tactic| llbc_verify_cond $f $dh) =>
    `(tactic|
      refine (_, (), ?_, rfl)
      simp only [evalFun, evalFunBody, evalStmtFuel, buildLocals,
        evalRValuePure, evalBinOpPure, evalOperandPure,
        evalPlacePure, writePlacePure, List.set, List.replicate,
        List.foldl, $f:ident, $dh:term])
\end{lstlisting}

Used for comparison-based functions (\texttt{max}, \texttt{min}, \texttt{abs}):

\begin{lstlisting}
theorem elab_max_ge (a b : Int) (h : a >= b) :
    evalFun [] maxFun 15 [.int a, .int b] at () |=
      .pureOutput (. = .int a) := by
  llbc_verify_cond maxFun (decide_eq_true h)
\end{lstlisting}

\subsubsection{llbc\_verify\_loop: Ground Loop Instances}

\begin{lstlisting}[caption={llbc\_verify\_loop macro.},label={lst:lvl}]
macro "llbc_verify_loop" v:term : tactic =>
  `(tactic| exact ($v:term, (), rfl, rfl))
\end{lstlisting}

Used for ground-instance loop theorems:

\begin{lstlisting}
theorem elab_sumTo_five : evalFun [] sumToFun 20 [.int 5] at () |=
    .pureOutput (. = .int 10) := by
  llbc_verify_loop (.int 10)
\end{lstlisting}

\subsection{Design Pitfalls and Lessons}

Writing Lean~4 macros involves non-obvious parsing constraints. We document
two issues encountered during development.

\subsubsection{Greedy Term Parsing}

The first attempt at \texttt{llbc\_verify\_prop} used two consecutive
\texttt{term} parameters:

\begin{lstlisting}[caption={Broken macro definition with two term params.},label={lst:broken}]
-- WRONG: Lean parses this as llbc_verify_prop (addFun h) ???
macro "llbc_verify_prop" f:term h:term : tactic => ...
\end{lstlisting}

Lean's parser is greedy: when it sees \texttt{llbc\_verify\_prop addFun h},
it reads \texttt{addFun} as the first term, then tries to extend it with
\texttt{h}, interpreting \texttt{addFun h} as a function application.
The second \texttt{term} parameter then has nothing to parse.

The fix is to make the first parameter an \texttt{ident}:

\begin{lstlisting}[caption={Fixed macro definition.},label={lst:fixed}]
-- CORRECT: ident is non-greedy, stops at whitespace
syntax "llbc_verify_prop" ident term : tactic
\end{lstlisting}

With \texttt{ident}, the parser correctly reads \texttt{addFun} as the
first parameter and \texttt{h} as the second.

\subsubsection{The \texttt{first |} Wrapper}

In \texttt{llbc\_verify}, the macro body initially used bare \texttt{refine}:

\begin{lstlisting}
-- WRONG: triggers "unexpected identifier; expected ')'" error
`(tactic| refine (_, (), ?_, rfl); simp only [...])
\end{lstlisting}

The \texttt{)} character after \texttt{rfl} confused the parser's
brace-matching in the macro body. The fix wraps the alternative in
\texttt{first |} to give the parser a clear delimiter:

\begin{lstlisting}
-- CORRECT:
`(tactic| exact (_, (), rfl, rfl) <|>
  (refine (_, (), ?_, rfl); simp only [...]))
\end{lstlisting}

These issues are not specific to this project; any Lean~4 tactic macro with
multiple tactic steps or nested term arguments is likely to encounter them.
We believe documenting them here benefits the Lean~4 proof engineering
community.

\subsection{Compression Measurement}

Table~\ref{tab:compression} compares manual proof size against the one-line
macro call for representative theorems.

\begin{table}[t]
\centering
\caption{Proof size comparison: manual \vs macro.}
\label{tab:compression}
\small
\begin{tabular}{@{}lrrrl@{}}
\toprule
\textbf{Theorem} & \textbf{Manual (lines)} & \textbf{Macro (lines)} & \textbf{Ratio} & \textbf{Macro used} \\
\midrule
\texttt{elab\_return42}   & 8  & 1 & 8$\times$  & \texttt{llbc\_verify} \\
\texttt{elab\_neg}        & 8  & 1 & 8$\times$  & \texttt{llbc\_verify} \\
\texttt{elab\_id}         & 8  & 1 & 8$\times$  & \texttt{llbc\_verify} \\
\texttt{elab\_add}        & 13 & 1 & 13$\times$ & \texttt{llbc\_verify\_prop} \\
\texttt{elab\_sub}        & 13 & 1 & 13$\times$ & \texttt{llbc\_verify\_prop} \\
\texttt{elab\_mul\_ok}    & 13 & 1 & 13$\times$ & \texttt{llbc\_verify\_prop} \\
\texttt{elab\_max\_ge}    & 12 & 1 & 12$\times$ & \texttt{llbc\_verify\_cond} \\
\texttt{elab\_min\_lt}    & 12 & 1 & 12$\times$ & \texttt{llbc\_verify\_cond} \\
\texttt{elab\_abs\_neg}   & 12 & 1 & 12$\times$ & \texttt{llbc\_verify\_cond} \\
\texttt{elab\_sumTo\_five}& 10 & 1 & 10$\times$ & \texttt{llbc\_verify\_loop} \\
\texttt{elab\_fib\_ten}   & 10 & 1 & 10$\times$ & \texttt{llbc\_verify\_loop} \\
\midrule
Average                    & 11.0 & 1 & 11$\times$ & \\
\bottomrule
\end{tabular}
\end{table}

The \texttt{Tactic/Examples.lean} module contains 16 theorems proved using
these macros. All 16 pass \texttt{lake build}, confirming that the macros
are syntactically correct and semantically sound. The average compression
ratio across the 16 theorems is approximately 11--12$\times$.

\subsection{Limits of the Automation}

The four macros handle the ``flat'' proof classes: pure reduction, overflow,
conditional, and ground loops. They do not automate loop invariant proofs.
The induction structure in \texttt{sumTo\_loop\_correct} is not uniform
across functions: \texttt{sumTo} uses an additive accumulator, \texttt{fact}
uses a multiplicative one, and \texttt{fib} uses two accumulators with a
Fibonacci recurrence. No single macro captures all three.

Automating loop invariant generation and proof is an active research
problem~\cite{verus2023,creusot2022}. Tools like Verus and Creusot use SMT
solvers to attempt automatic invariant synthesis. Our approach is complementary:
we provide a clear, human-readable inductive structure that can be checked by
the Lean kernel, at the cost of requiring the invariant to be supplied
manually.

%% ============================================================
%%  §8  TypeScript Extension and Cross-Language Unification
%% ============================================================

\section{TypeScript Extension and Cross-Language Unification}
\label{sec:typescript}

The TypeScript extension demonstrates that the \RustM monad and \ProgramSpec
specification language are genuinely language-agnostic.
\texttt{TypeScript/ElabSpec.lean} contains \textbf{33 theorems} (0 sorry)
covering 14 TypeScript functions — including a \texttt{while} loop — and
6 cross-language unification theorems that formally assert both a Rust and a
TypeScript implementation satisfy the \emph{same} specification.
An additional \textbf{6 theorems} in \texttt{TypeScript/BugDetection.lean}
demonstrate how the kernel acts as a bug oracle: it catches two classic
programming errors and confirms the fixes, all by \rfl.

\subsection{The Verified TypeScript Fragment}
\label{sec:ts-fragment}

TypeScript's execution model differs substantially from Rust's: dynamic
typing, garbage collection, prototype-based objects, no ownership.
The fragment we verify is a well-defined subset:

\begin{center}
\small
\begin{tabular}{@{}ll@{}}
\toprule
\textbf{Feature}                  & \textbf{Status} \\
\midrule
\texttt{number} (modelled as $\mathbb{Z}$) & Supported \\
\texttt{boolean}                   & Supported \\
\texttt{string} (constants and concatenation) & Partial (in evaluator; no theorems) \\
Arithmetic, comparison, logical ops & Supported \\
Strict equality (\texttt{===}), strict inequality (\texttt{!==}) & Supported \\
\texttt{if}/\texttt{else}, ternary \texttt{?:} & Supported \\
\texttt{while} loops with mutable variable update & Supported \\
\texttt{let}/\texttt{const} bindings & Supported \\
\texttt{typeof}                    & Supported \\
\midrule
Closures, arrow functions          & \textbf{Not supported} \\
Objects (\texttt{\{\}}), prototype chain & \textbf{Not supported} \\
Arrays (indexing in evaluator; no theorems) & Partial \\
\texttt{async}/\texttt{await}, Promises & \textbf{Not supported} \\
Type coercions (\texttt{==}, \texttt{+} on mixed types) & \textbf{Not supported} \\
IEEE~754 floating-point (NaN, Infinity, \texttt{-0}) & \textbf{Not supported} \\
\bottomrule
\end{tabular}
\end{center}

\noindent
The \texttt{number} type is modelled as $\mathbb{Z}$ (Lean's \texttt{Int})
rather than IEEE~754 \texttt{double}. This is a deliberate simplification:
integer arithmetic is exact under this model, which is sound for programs
that only use integer-valued expressions. Programs depending on floating-point
behaviour (NaN propagation, rounding, \texttt{-0}) are outside scope.

\subsection{The TypeScript AST and Evaluator}

\begin{lstlisting}[caption={TypeScript AST (actual definitions, selected).},
                   label={lst:tsast}]
inductive TSExpr where
  | lit   (l : TSLit)                            -- numeric/boolean literal
  | var   (idx : Nat)                            -- positional; idx 0 = first param
  | binOp (op : String) (l r : TSExpr)
  | unOp  (op : String) (e : TSExpr)
  | ite   (cond thenE elseE : TSExpr)

inductive TSStmt where
  | skip
  | return_ (e : TSExpr)
  | const   (name : String) (ty : Option TSTy) (e : TSExpr) (k : TSStmt)
  | set_    (idx : Nat) (e : TSExpr) (k : TSStmt)  -- mutable update env[idx]
  | ite     (cond : TSExpr) (thenB elseB : TSStmt)
  | while_  (cond : TSExpr) (body : TSStmt)
  | seq     (s1 s2 : TSStmt)
\end{lstlisting}

\noindent
\textbf{Key design choices:}
\begin{itemize}[leftmargin=2em]
  \item \textbf{Positional environments.} Variables are de Bruijn indices
        (\texttt{var idx}) into a \texttt{List TSValue}. This makes
        \texttt{evalExpr env (.var n) = env[n]?} kernel-transparent without
        any string equality, enabling all branching proofs to go through by
        \texttt{decide\_eq\_true/false} without sorry.

  \item \textbf{\texttt{const} + \texttt{set\_}.} Immutable bindings push a
        value to the front of the environment; \texttt{set\_} updates an existing
        slot in place. This pair covers both functional-style local declarations
        and imperative loop-variable mutation.

  \item \textbf{Same monad.} \texttt{evalTSFun} returns \texttt{RustM Unit TSValue}
        — exactly the type returned by \texttt{evalFun} for Rust. The common
        return type is what makes cross-language unification statements type-check.
\end{itemize}

The evaluator follows the same fuel-parametric, \texttt{match}-on-\texttt{Except}
discipline as the LLBC interpreter:

\begin{lstlisting}[caption={TypeScript evaluator entry point.},label={lst:tseval}]
def evalTSFun (f : TSFunDef) (fuel : Nat) (args : List TSValue) :
    RustM Unit TSValue :=
  fun _ =>
    let initEnv : Env := args     -- positional; idx 0 = first arg
    match evalStmt fuel f.body initEnv with
    | Except.ok (.return_ v, _) => Except.ok (v, ())
    | Except.ok (.next,     _)  => Except.error (.explicit "fell off end")
    | Except.error e            => Except.error e
    | _                         => Except.error (.explicit "break outside loop")
\end{lstlisting}

\subsection{Verified TypeScript Functions}

The 14 functions proved correct are:
\texttt{return42}, \texttt{id}, \texttt{neg}, \texttt{add}, \texttt{sub},
\texttt{max}, \texttt{isZero}, \texttt{mul}, \texttt{min},
\texttt{abs}, \texttt{not\_gate}, \texttt{clamp}, and \texttt{sum\_to}.
The first 12 are pure functions; \texttt{sum\_to} uses a \texttt{while} loop.

\paragraph{Pure functions.}
For functions with no branches (e.g., \texttt{add}, \texttt{neg}), the proof
is by \rfl — the kernel reduces the closed term to its value:

\begin{lstlisting}[caption={T3: neg proved by rfl.},label={lst:tsneg}]
def tsNegFun : TSFunDef := {
  body := .return_ (.unOp "-" (.var 0))  -- x = args[0]
  ...
}
theorem elab_ts_neg (x : Int) :
    evalTSFun tsNegFun 10 [.num x] at () |= .pureOutput (· = .num (-x)) :=
  ⟨.num (-x), (), rfl, rfl⟩
\end{lstlisting}

\paragraph{Branching functions.}
For functions with conditionals (\texttt{max}, \texttt{abs}, \texttt{clamp}),
we supply a branch hypothesis and use \texttt{decide\_eq\_true/false}
in a two-pass simp:

\begin{lstlisting}[caption={T10: abs (x $\ge$ 0 branch).},label={lst:tsabs}]
theorem elab_ts_abs_nonneg (x : Int) (h : x ≥ 0) :
    evalTSFun tsAbsFun 10 [.num x] at () |= .pureOutput (· = .num x) := by
  refine ⟨.num x, (), ?_, rfl⟩
  simp only [evalTSFun, tsAbsFun]
  simp only [evalStmt, evalExpr, evalBinOp, evalLit,
             List.getElem?_cons_zero, decide_eq_true h]
\end{lstlisting}

\paragraph{The while loop: sum\_to.}
\texttt{sum\_to(n)} computes $0 + 1 + \cdots + (n-1)$.
In TypeScript, this is a \texttt{while} loop over two mutable variables
\texttt{s} (accumulator) and \texttt{i} (counter):

\begin{lstlisting}[language=TypeScript,caption={TypeScript sum\_to source.},label={lst:tssumto}]
function sum_to(n: number): number {
  let s = 0;
  let i = 0;
  while (i < n) { s += i; i += 1; }
  return s;
}
\end{lstlisting}

The \texttt{TSFunDef} uses \texttt{const} to push \texttt{s} and \texttt{i}
onto the positional environment (making \texttt{i = env[0]},
\texttt{s = env[1]}, \texttt{n = env[2]}), and \texttt{set\_} to update
them in-place each iteration:

\begin{lstlisting}[caption={T14: tsSumToFun definition.},label={lst:tssumtofun}]
def tsSumToFun : TSFunDef := {
  name := "sum_to"; params := [⟨"n", .number, false⟩]; retTy := .number
  body :=
    .const "s" none (.lit (.num 0))     -- env = [0, n]
    (.const "i" none (.lit (.num 0))    -- env = [0, 0, n]; i=0,s=1,n=2
    (.seq
      (.while_ (.binOp "<" (.var 0) (.var 2))   -- while i < n
        (.set_ 1 (.binOp "+" (.var 1) (.var 0)) -- s = s + i
        (.set_ 0 (.binOp "+" (.var 0) (.lit (.num 1))) -- i = i + 1
          .skip)))
      (.return_ (.var 1))))             -- return s
  async_ := false
}
\end{lstlisting}

The correctness theorems are ground instances proved by \rfl (the kernel
evaluates the loop concretely):

\begin{lstlisting}[caption={T14 ground instances.},label={lst:tssumtothm}]
theorem elab_ts_sumTo_five :
    evalTSFun tsSumToFun 100 [.num 5] at () |= .pureOutput (· = .num 10) :=
  ⟨.num 10, (), rfl, rfl⟩

theorem elab_ts_sumTo_ten :
    evalTSFun tsSumToFun 100 [.num 10] at () |= .pureOutput (· = .num 45) :=
  ⟨.num 45, (), rfl, rfl⟩
\end{lstlisting}

\subsection{Cross-Language Unification Theorems (U1--U6)}
\label{sec:unification}

Six unification theorems establish that Rust and TypeScript implementations
satisfy the \emph{same} \ProgramSpec predicate:

\begin{lstlisting}[caption={U1: max — same spec from both languages.},
                   label={lst:unified}]
theorem unified_max_either (a b : Int) :
    (evalFun [] maxFun 10 [.int a, .int b] at () |=
      .pureOutput (fun v => v = .int a ∨ v = .int b)) ∧
    (evalTSFun tsMaxFun 10 [.num a, .num b] at () |=
      .pureOutput (fun v => v = .num a ∨ v = .num b)) :=
  ⟨elab_max_either a b, elab_ts_max_either a b⟩
\end{lstlisting}

\begin{lstlisting}[caption={U5: sum\_to — same spec from Rust loop and TypeScript while loop.},
                   label={lst:unifiedsumto}]
theorem unified_sumTo_ten :
    (evalFun [] sumToFun 1000 [.int 10] at () |=
      .pureOutput (· = .int 45)) ∧
    (evalTSFun tsSumToFun 100 [.num 10] at () |=
      .pureOutput (· = .num 45)) :=
  ⟨elab_sumTo_ten, elab_ts_sumTo_ten⟩
\end{lstlisting}

\begin{lstlisting}[caption={U6: neg — symbolic unification.},label={lst:unifiedneg}]
theorem unified_neg (x : Int) :
    (evalFun [] negFun 10 [.int x] at () |= .pureOutput (· = .int (-x))) ∧
    (evalTSFun tsNegFun 10 [.num x] at () |= .pureOutput (· = .num (-x))) :=
  ⟨elab_neg x, elab_ts_neg x⟩
\end{lstlisting}

\noindent
The full set of 6 unification theorems (U1--U6) covers
\texttt{max} (either branch), \texttt{add} (no-crash), \texttt{min} (either),
\texttt{mul} (no-crash), \texttt{sum\_to} (ground, n=10), and \texttt{neg} (symbolic).

These theorems are the formal realisation of the paper's core claim:
a single \ProgramSpec instance describes implementations in both languages.
The proofs are simple conjunctions of the individual results; the substantive
work is in establishing compatible specifications on each side.
From a software-engineering perspective, if a function is later ported from
Rust to TypeScript (or vice versa), the unification theorem formally certifies
that the rewrite preserves the specification.

\subsection{Bug Detection and Repair}
\label{sec:ts-bug-detection}

A key practical benefit of the deep-embedding approach is that the Lean kernel
acts as a \emph{bug oracle}: if a function definition contains an error, a proof
attempt immediately reveals the wrong computed value through a failed \rfl reduction.
\texttt{TypeScript/BugDetection.lean} demonstrates two representative bugs.

\paragraph{Bug B1: off-by-one in \texttt{sum\_to}.}
The intended function computes $0 + 1 + \cdots + (n-1)$.
The buggy version uses \texttt{<=} instead of \texttt{<} in the loop guard:

\begin{lstlisting}[language=TypeScript,caption={Buggy TypeScript sum\_to.},label={lst:buggysum}]
// BUG: condition should be i < n
while (i <= n) { s += i; i += 1; }
\end{lstlisting}

\noindent
This causes the loop to run for $i = 0, 1, \ldots, n$ (inclusive), computing
$n(n{+}1)/2$ instead of $n(n{-}1)/2$.
Attempting to prove \texttt{sum\_to\_buggy(5) = 10} fails:
\texttt{⟨.num 10, (), rfl, rfl⟩} does not type-check because the kernel
reduces the term to \texttt{.num 15}.
Theorem B1 instead \emph{documents the bug} by proving the wrong value:

\begin{lstlisting}[caption={B1: kernel confirms the off-by-one bug.},label={lst:b1}]
theorem ts_sumTo_buggy_five :
    evalTSFun tsSumToBuggyFun 100 [.num 5] at () |= .pureOutput (· = .num 15) :=
  ⟨.num 15, (), rfl, rfl⟩   -- 15 = 0+1+2+3+4+5 (includes 5; should be 10)
\end{lstlisting}

\noindent
Fixing the guard to \texttt{<} recovers the correct \texttt{tsSumToFun}, for which
\texttt{elab\_ts\_sumTo\_five} (T13) proves the result is \texttt{.num 10}.

\paragraph{Bug B2: wrong zero guard in \texttt{safeDiv}.}
A safe-division function should return $a/b$ when $b \neq 0$ and $0$
otherwise.  The buggy version guards on $b > 0$, silently returning $0$
for negative divisors:

\begin{lstlisting}[language=TypeScript,caption={Buggy TypeScript safeDiv.},label={lst:buggydiv}]
// BUG: guard should be b !== 0
function safeDivBuggy(a: number, b: number): number {
  if (b > 0) return a / b;
  return 0;
}
\end{lstlisting}

\noindent
Theorem B3 exposes the bug: the kernel evaluates \texttt{safeDivBuggy(10,~-2)}
to \texttt{0}, not \texttt{-5}:

\begin{lstlisting}[caption={B3: kernel exposes the wrong-guard bug.},label={lst:b3}]
theorem ts_safeDiv_buggy_negative_divisor :
    evalTSFun tsSafeDivBuggyFun 10 [.num 10, .num (-2)] at ()
      |= .pureOutput (· = .num 0) :=
  ⟨.num 0, (), rfl, rfl⟩    -- returns 0 for b=-2 (bug!)
\end{lstlisting}

\noindent
Fixing the guard to \texttt{!==} yields \texttt{tsSafeDivFun}.
Theorems B4--B6 verify all three cases: negative divisor returns $-5$,
positive divisor returns $3$, and zero divisor safely returns $0$
without panicking:

\begin{lstlisting}[caption={B4--B6: fixed safeDiv verified for all divisor signs.},label={lst:b46}]
theorem ts_safeDiv_correct_neg :
    evalTSFun tsSafeDivFun 10 [.num 10, .num (-2)] at () |= .pureOutput (· = .num (-5)) :=
  ⟨.num (-5), (), rfl, rfl⟩

theorem ts_safeDiv_correct_zero :
    evalTSFun tsSafeDivFun 10 [.num 10, .num 0] at () |= .pureOutput (· = .num 0) :=
  ⟨.num 0, (), rfl, rfl⟩
\end{lstlisting}

\noindent
All six bug-detection theorems are proved by \rfl.
The proof methodology is the same as for correctness proofs: the only
difference is that for the buggy theorems, the stated output value is the
\emph{wrong} answer, serving as machine-certified documentation of the defect.

\subsection{Absence of Sorry}

All 39 TypeScript theorems (33 in \texttt{ElabSpec.lean} + 6 in \texttt{BugDetection.lean})
are fully kernel-checked (0 sorry).
The key enabler is the \textbf{positional environment}: because
\texttt{env[n]?} reduces by \texttt{List.getElem?} without any
\texttt{@[extern]} string-equality instance, all evaluation equations
reduce under \texttt{simp}. Branching theorems supply a hypothesis
\texttt{h} and discharge the \texttt{decide} call via
\texttt{decide\_eq\_true h} or \texttt{decide\_eq\_false (not\_le.mpr h)}.

%% ============================================================
%%  §9  Fibonacci Case Study
%% ============================================================

\section{The Fibonacci Case Study}
\label{sec:fib}

The Fibonacci function is a natural benchmark for loop invariant methodology:
it maintains two accumulator variables simultaneously, and the invariant
involves a recurrence relation (\texttt{Nat.fib\_add\_two}) rather than
a simple arithmetic formula. This section presents the complete zero-\sorry
proof of \texttt{fib\_correct}.

\subsection{The Function}

The Rust implementation of Fibonacci uses the two-accumulator pattern:
start with $(a, b) = (0, 1)$ representing $(F(0), F(1))$, and in each
iteration transform $(a, b) \mapsto (b, a+b)$, advancing the index by one.
After $n$ iterations, $a = F(n)$ and $b = F(n+1)$.

The LLBC representation allocates 8 local variables:
\begin{center}
\texttt{[0]=ret, [1]=n, [2]=a, [3]=b, [4]=i, [5]=t, [6]=cond, [7]=tmp}
\end{center}

One loop iteration executes:
\begin{enumerate}[leftmargin=2em]
  \item \texttt{t   := b}         (save old $b$)
  \item \texttt{tmp := a + b}     (compute $a + b$)
  \item \texttt{b   := tmp}       ($b \leftarrow a + b$)
  \item \texttt{a   := t}         ($a \leftarrow$ old $b$)
  \item \texttt{tmp := i + 1}
  \item \texttt{i   := tmp}       ($i \leftarrow i + 1$)
\end{enumerate}

The body (in LLBC) therefore requires 9 fuel units in the true branch
(6 assignments + 2 for condition + \texttt{ite}) and 3 units in the false
branch (condition + \texttt{ite} + \texttt{break\_}).

\subsection{The Invariant Structure}

\begin{definition}[Two-accumulator Fibonacci invariant]
After $k$ iterations starting from state $(a, b, i)$ where
$a = F(i)$ and $b = F(i+1)$, the loop terminates with
$(a', b', i') = (F(n), F(n+1), n)$.
\end{definition}

This is stronger than the typical single-accumulator invariant: we must
track both $a$ and $b$ simultaneously, because neither alone carries enough
information to reconstruct the Fibonacci sequence.

The overflow predicate is:

\begin{lstlisting}[caption={FibOv: overflow condition for fib.},label={lst:fibov}]
def FibOv (n i : Int) : Prop :=
  forall  j : Int, i <= j -> j < n ->
    (minI32 <= ((Nat.fib j.toNat) : Int) + (Nat.fib (j.toNat + 1)) /\
     ((Nat.fib j.toNat) : Int) + (Nat.fib (j.toNat + 1)) <= maxI32) /\
    (minI32 <= j + 1 /\ j + 1 <= maxI32)
\end{lstlisting}

Note that \FibOv is parametric in $i$ (the current loop index) but not
in $a$ and $b$ directly — because the invariant pins $a = F(i)$ and $b = F(i+1)$,
the overflow condition can be expressed entirely in terms of Fibonacci values.
This makes \FibOv independent of the current accumulator values, simplifying
the induction.

\subsection{Step Lemmas}

Two step lemmas set up the induction.

\begin{lstlisting}[caption={fib\_body\_true: true branch step lemma.},label={lst:fibbt}]
theorem fib_body_true (env) (n a b i : Int) (c5 c6 c7 : Value)
    (h    : i < n)
    (hov1 : minI32 <= a + b /\ a + b <= maxI32)
    (hov2 : minI32 <= i + 1 /\ i + 1 <= maxI32)
    (fuel : Nat) (hf : 9 <= fuel) :
    evalStmtFuel env fuel fibBody
      [.unit, .int n, .int a, .int b, .int i, c5, c6, c7] =
    .ok (.next, [.unit, .int n, .int b, .int (a + b), .int (i + 1),
                 .int b, .bool_ true, .int (i + 1)]) := by
  obtain (f, rfl) := ...
  have hdt : decide (i < n) = true := decide_eq_true h
  simp only [fibBody, evalStmtFuel, ..., hdt, if_pos hov1, if_pos hov2]
\end{lstlisting}

\begin{lstlisting}[caption={fib\_body\_false: false branch step lemma.},label={lst:fibbf}]
theorem fib_body_false (env) (n a b i : Int) (c5 c6 c7 : Value)
    (h    : ¬(i < n))
    (fuel : Nat) (hf : 3 <= fuel) :
    evalStmtFuel env fuel fibBody
      [.unit, .int n, .int a, .int b, .int i, c5, c6, c7] =
    .ok (.break_, [.unit, .int n, .int a, .int b, .int i,
                   c5, .bool_ false, c7]) := by
  obtain (f, rfl) := ...
  have hdf : decide (i < n) = false := decide_eq_false h
  simp only [fibBody, evalStmtFuel, ..., hdf]
\end{lstlisting}

The overflow hypothesis \texttt{hov1} in \texttt{fib\_body\_true} is required
because the \texttt{tmp := a + b} assignment calls \texttt{evalBinOpPure .add},
which checks $\texttt{minI32} \le a + b \le \texttt{maxI32}$. The \texttt{hov2}
hypothesis covers the \texttt{tmp := i + 1} step.

\subsection{The Loop-Correct Theorem}

\begin{lstlisting}[caption={fib\_loop\_correct: inductive invariant.},label={lst:fiblc}]
theorem fib_loop_correct (k : Nat) (env : FunEnv)
    (n a b i : Int) (c5 c6 c7 : Value)
    (heq   : (n - i).toNat = k)
    (hi    : 0 <= i)  (hin : i <= n)
    (hinvA : a = (Nat.fib i.toNat))
    (hinvB : b = (Nat.fib (i.toNat + 1)))
    (hov   : FibOv n i)
    (fuel  : Nat) (hfuel : k + 10 <= fuel) :
    exists  s5 s6 s7 : Value,
    evalStmtFuel env fuel (.loop fibBody)
      [.unit, .int n, .int a, .int b, .int i, c5, c6, c7] =
    .ok (.next, [.unit, .int n, .int (Nat.fib n.toNat),
                 .int (Nat.fib (n.toNat + 1)), .int n, s5, s6, s7])
\end{lstlisting}

The proof proceeds by induction on $k = (n - i).\texttt{toNat}$:

\textbf{Zero case ($k = 0$, \ie $i = n$).}
From $i = n$ and the invariant hypotheses, we have $a = F(n)$ and
$b = F(n+1)$. The loop body evaluates the false branch (\texttt{fib\_body\_false}),
exiting immediately. The goal closes by \texttt{simp only [evalStmtFuel, hbody]}.

\textbf{Successor case ($k = k'+1$, \ie $i < n$).}
First, from $i < n$ and $\FibOv\,n\,i$, we extract the overflow bounds at $j = i$
using \texttt{FibOv\_self}. Since $a = F(i)$ and $b = F(i+1)$ (from the
invariant hypotheses), the body executes the true branch.
After one iteration, the new state is $(b, a+b, i+1)$. We establish:

\begin{itemize}[leftmargin=1.5em]
  \item $b = F(i+1) = F((i+1).\texttt{toNat})$ (from \texttt{hinvB} and the
        fact that $(i+1).\texttt{toNat} = i.\texttt{toNat} + 1$).
  \item $a + b = F(i) + F(i+1) = F(i+2) = F((i+1).\texttt{toNat} + 1)$,
        using \texttt{Nat.fib\_add\_two}.
\end{itemize}

The induction hypothesis then applies with $k'$ remaining iterations,
proving the result.

The key Mathlib lemma is:

\begin{lstlisting}[caption={Nat.fib\_add\_two from Mathlib.},label={lst:fibrecurrence}]
-- From Mathlib: Fibonacci recurrence relation
theorem Nat.fib_add_two (n : Nat) :
    Nat.fib (n + 2) = Nat.fib n + Nat.fib (n + 1) := ...
\end{lstlisting}

In the successor case we use it in the form $F(i) + F(i+1) = F(i+2)$:

\begin{lstlisting}[caption={Using fib\_add\_two in the invariant step.},label={lst:fibstep}]
have hinvB' : a + b = (Nat.fib ((i + 1).toNat + 1)) := by
  rw [hinvA, hinvB, hi1,
      show i.toNat + 1 + 1 = i.toNat + 2 from by omega,
      <- Nat.cast_add, <- Nat.fib_add_two]
\end{lstlisting}

\subsection{Safety Condition and native\_decide}

\begin{definition}[FibSafe]
$\texttt{FibSafe}(n) \;\Leftrightarrow\; 0 \le n \le 45$.
\end{definition}

The choice of 45 as the upper bound comes from the value of $F(46)$:

\begin{lstlisting}[caption={Fibonacci bound: F(46) fits in I32.},label={lst:f46}]
-- F(46) = 1836311903 < 2147483647 = maxI32
-- F(47) = 2971215073 > 2147483647 = maxI32
-- So fib is safe exactly for n <= 45.
have hfib46 : Nat.fib 46 = 1836311903 := by native_decide
\end{lstlisting}

The \texttt{native\_decide} tactic runs Lean's native code compiler to evaluate
\texttt{Nat.fib 46} and confirm that it equals 1\,836\,311\,903. This is
\emph{not} \sorry: \texttt{native\_decide} generates a verified proof term
by running the Lean kernel's native evaluator. It is appropriate here because
computing $F(46)$ by kernel reduction would be prohibitively slow (Fibonacci
grows exponentially), but native evaluation is fast.

Using this fact, \texttt{fibSafe\_implies\_ov} establishes \FibOv for any
$n \le 45$:

\begin{lstlisting}[caption={fibSafe\_implies\_ov proof structure.},label={lst:fibsafeov}]
theorem fibSafe_implies_ov (n : Int) (hsafe : FibSafe n) :
    FibOv n 0 := by
  intro j hj hjn
  have hjnat2 : j.toNat + 2 <= 46 := by omega
  refine ((?_, ?_), ?_)
  . -- lower bound: Fib >= 0
    have : (0 : Int) <= (Nat.fib j.toNat) := Int.natCast_nonneg _
    norm_num [IntBounds.minI32]; linarith
  . -- upper bound: F(j+2) <= F(46) = 1836311903 <= maxI32
    have hstep : Nat.fib j.toNat + Nat.fib (j.toNat + 1) =
                 Nat.fib (j.toNat + 2) := by rw [<- Nat.fib_add_two]
    have hmono : Nat.fib (j.toNat + 2) <= Nat.fib 46 :=
                 Nat.fib_mono hjnat2
    have hfib46 : Nat.fib 46 = 1836311903 := by native_decide
    have hle_nat : Nat.fib (j.toNat + 2) <= 1836311903 := hfib46 ▸ hmono
    have hle_int : (Nat.fib (j.toNat + 2) : Int) <= 1836311903 :=
                   by exact_mod_cast hle_nat
    have hsum : ... := by exact_mod_cast hstep
    norm_num [IntBounds.maxI32]; linarith
  . -- j+1 bounds: j+1 <= 45 <= maxI32
    exact (by norm_num [IntBounds.minI32]; linarith,
           by norm_num [IntBounds.maxI32]; linarith)
\end{lstlisting}

\subsection{The Main Theorem}

\begin{lstlisting}[caption={fib\_correct: the main theorem.},label={lst:fibcorrect}]
def FibSafe (n : Int) : Prop := 0 <= n /\ n <= 45

theorem fib_correct (n : Int) (hsafe : FibSafe n) :
    evalFun [] fibFun (n.toNat + 14) [.int n] at () |=
      .pureOutput (. = .int (Nat.fib n.toNat)) := by
  have hov := fibSafe_implies_ov n hsafe
  obtain (hn, _) := hsafe
  -- Initial invariant: a = F(0) = 0, b = F(1) = 1
  obtain (s5, s6, s7, hloop) :=
    fib_loop_correct n.toNat (fibFun :: []) n 0 1 0 .unit .unit .unit
      (by omega) (by omega) hn
      (by simp [Nat.fib]) (by simp [Nat.fib])
      hov (n.toNat + 10) (by omega)
  refine (.int (Nat.fib n.toNat), (), ?_, rfl)
  -- hpre reduces the 3 prefix assigns + seq without unfolding the loop
  have hpre : ... := rfl
  simp only [evalFun, evalFunBody, hpre, hloop]
  simp only [evalStmtFuel, evalRValuePure, evalOperandPure,
    evalPlacePure, writePlacePure, List.set, ...]
\end{lstlisting}

The fuel bound \texttt{n.toNat + 14} decomposes as follows:
\begin{itemize}[leftmargin=2em]
  \item 3 prefix assignments (initialising \texttt{a}, \texttt{b}, \texttt{i}): 3 fuel
  \item 1 \texttt{seq} between the prefix and the loop: 1 fuel
  \item Loop body: $k + 10$ fuel where $k = n.\texttt{toNat}$ and 10 is the
        per-iteration overhead (9 for the true branch + 1 for the outer loop step)
\end{itemize}

Total: $3 + 1 + n.\texttt{toNat} + 10 = n.\texttt{toNat} + 14$.

\subsection{Summary}

The Fibonacci proof consists of 9 theorems with 0 \sorry:
\texttt{fib\_body\_true}, \texttt{fib\_body\_false}, \texttt{FibOv\_self},
\texttt{FibOv\_step}, \texttt{fib\_loop\_correct},
\texttt{fibSafe\_implies\_ov}, \texttt{fib\_correct}, plus two auxiliary
\texttt{getElem?} lemmas. The use of \texttt{native\_decide} for
$F(46) = 1\,836\,311\,903$ is the only instance of non-kernel reasoning in
the entire file; all other steps are kernel-checked. The proof is approximately
230 lines of Lean~4 code, including documentation comments.

Compared to the \texttt{sumTo} and \texttt{fact} proofs, the key novel
element is the two-accumulator invariant: both \texttt{hinvA} and
\texttt{hinvB} must be maintained and updated in the inductive step, and
the update of \texttt{hinvB} requires the Fibonacci recurrence relation
from Mathlib. This pattern extends naturally to other linear recurrences
(Lucas numbers, tribonacci) by replacing \texttt{Nat.fib\_add\_two} with
the appropriate recurrence theorem.

%% ============================================================
%%  §10  Evaluation
%% ============================================================

\section{Evaluation}
\label{sec:evaluation}

This section reports quantitative measurements of the lean-pl-verify artifact:
theorem counts, lines of code, the end-to-end Charon pipeline,
proof-to-code ratios, build times, and a comparison with related work.
The complete project builds with \textbf{0 sorry} across all 258 theorems.

\subsection{Theorem Counts}

Table~\ref{tab:theorems} gives the complete theorem count by module.
All theorems are fully kernel-checked.
\texttt{Translation/Adequacy.lean} proves both directions of adequacy:
soundness (every interpreter output is relationally derivable) and
completeness (every relational derivation is witnessed by some fuel bound).

\begin{table}[t]
\centering
\caption{Theorem counts by module. Sorry count: \textbf{0} across all files.}
\label{tab:theorems}
\small
\begin{tabular}{@{}llrr@{}}
\toprule
\textbf{Module} & \textbf{Contents} & \textbf{Theorems} & \textbf{Sorry} \\
\midrule
\multicolumn{4}{@{}l}{\textit{Foundation}} \\
\texttt{Foundation/Monad.lean}        & Monad laws, rpanic lemmas              &  2 & 0 \\
\texttt{Foundation/Ownership.lean}    & FBPair identity laws                   &  4 & 0 \\
\midrule
\multicolumn{4}{@{}l}{\textit{Specification}} \\
\texttt{Spec/Satisfies.lean}          & ProgramSpec combinators                & 17 & 0 \\
\texttt{Spec/Examples.lean}           & RustM-level examples (E1--E8)          & 15 & 0 \\
\midrule
\multicolumn{4}{@{}l}{\textit{LLBC / Rust verification}} \\
\texttt{Translation/ElabSpec.lean}    & LLBC Rust theorems (F1--F16)           & 48 & 0 \\
\texttt{Translation/LoopInvariant.lean} & \texttt{sum\_to} loop invariant      & 14 & 0 \\
\texttt{Translation/FactInvariant.lean} & \texttt{fact} loop invariant         & 14 & 0 \\
\texttt{Translation/FibInvariant.lean}  & \texttt{fib} loop invariant          &  9 & 0 \\
\texttt{Translation/Semantics.lean}   & Relational big-step semantics          &  0$^\dagger$ & 0 \\
\texttt{Translation/Adequacy.lean}    & Full adequacy (soundness + completeness) & 18 & 0 \\
\midrule
\multicolumn{4}{@{}l}{\textit{Charon-extracted Rust (end-to-end pipeline)}} \\
\texttt{Translation/CharonDefs.lean}  & Auto-generated \texttt{LLBCFunDef}s    &  0$^\ddagger$ & 0 \\
\texttt{Translation/CharonSpec.lean}  & Proofs over Charon-generated defs      & 57 & 0 \\
\midrule
\multicolumn{4}{@{}l}{\textit{Real-crate extraction (\texttt{num-integer} v0.1.45)}} \\
\texttt{Translation/NumIntegerSpec.lean} & \texttt{Integer::is\_even<u32>} from published crate & 5 & 0 \\
\midrule
\multicolumn{4}{@{}l}{\textit{Proof automation}} \\
\texttt{Tactic/Examples.lean}         & One-liner proofs via macros            & 16 & 0 \\
\midrule
\multicolumn{4}{@{}l}{\textit{TypeScript (cross-language unification)}} \\
\texttt{TypeScript/ElabSpec.lean}     & TS theorems T1--T14 + 6 unif.\ thms   & 33 & 0 \\
\texttt{TypeScript/BugDetection.lean} & Bug detection case study (B1--B6)      &  6 & 0 \\
\midrule
\textbf{Total}                        &                                        & \textbf{258} & \textbf{0} \\
\bottomrule
\end{tabular}
\end{table}

\noindent
$^\dagger$ \texttt{Semantics.lean} defines the relational big-step semantics
and supporting lemmas; its correctness is witnessed by the adequacy theorems
in \texttt{Adequacy.lean}.

\noindent
$^\ddagger$ \texttt{CharonDefs.lean} contains 18 auto-generated
\texttt{LLBCFunDef} values (definitions, not theorems); all 57
correctness theorems reside in \texttt{CharonSpec.lean}.

\subsection{The Charon End-to-End Pipeline}
\label{sec:eval-charon}

A central engineering contribution is the pipeline that connects Rust
source code to kernel-checked Lean proofs with no manual transcription:

\begin{enumerate}[leftmargin=2em]
  \item \textbf{Rust source} (\texttt{examples/rust-crate/src/lib.rs},
        148 lines, 18 functions).
  \item \textbf{Charon} v0.1.197 (\texttt{charon cargo --dest-file ...})
        translates MIR to LLBC JSON (\texttt{verified\_fns.llbc}).
  \item \textbf{charon2lean.py} (463 lines) translates the LLBC JSON to
        \texttt{LLBCFunDef} Lean terms, handling:
        flat-list-to-CPS conversion,
        \texttt{AddChecked}/\texttt{SubChecked}/\texttt{MulChecked} desugaring,
        overflow-check \texttt{Assert} elision,
        negative integer literal wrapping,
        and dict-form binary operators (\eg \texttt{\{"Rem": "UB"\}}).
  \item \textbf{\texttt{CharonDefs.lean}} (auto-generated) is
        elaborated by Lean's kernel without error.
  \item \textbf{\texttt{CharonSpec.lean}} proves
        57 correctness theorems over the auto-generated definitions
        using the same proof tactics as \texttt{ElabSpec.lean}.
\end{enumerate}

\paragraph{Pipeline validation: adding \texttt{pow} and \texttt{gcd}.}
To validate that the pipeline is extensible beyond the initial 16 functions,
we added \texttt{pow} (integer power, loop) and \texttt{gcd}
(Euclidean algorithm using \texttt{\%}).
This required: (1)~writing the Rust functions and re-running Charon;
(2)~fixing \texttt{charon2lean.py} to handle the \texttt{\{"Rem": "UB"\}}
dict-form operator that Charon emits for signed remainder;
(3)~adding \texttt{.bitAnd} on \texttt{bool\_} to \texttt{evalBinOpPure}
(Charon's UB-check intermediate uses bool-typed bitwise-and);
(4)~writing 10 new theorems.
The proofs are all by \texttt{rfl}: \texttt{gcd(a, 0) = a} and
\texttt{pow(x, 0) = 1} hold symbolically (loop exits immediately),
and ground instances (\texttt{gcd(12, 8) = 4}, \texttt{pow(2, 10) = 1024})
reduce by direct kernel evaluation.

The pipeline demonstrates that the proof infrastructure is \emph{not tied to
hand-crafted definitions}: Charon-generated ASTs, including Charon's
UB-check temporaries, are handled uniformly by \texttt{charon\_simp}.

\subsection{Real-Crate Extraction: \texttt{num-integer} v0.1.45}
\label{sec:eval-numinteger}

To demonstrate that the framework connects to real published crates, we
ran Charon directly on the source of \texttt{num-integer} v0.1.45
(\url{https://crates.io/crates/num-integer}, $\approx$45M downloads).

\paragraph{Extraction.}
Running \texttt{charon cargo} on the crate source produced an LLBC file
with 1\,657 function entries: 359 with concrete bodies, 590 opaque
(standard library intrinsics and trait stubs).
Among the concrete bodies is the macro-generated implementation of
\texttt{Integer::is\_even} for each primitive type.
We extracted the \texttt{u32} implementation (11 LLBC statements):

\begin{lstlisting}[caption={Source: num-integer/src/lib.rs (impl\_integer\_for\_usize! macro).},
                   label={lst:iseven}]
fn is_even(&self) -> bool { *self % 2 == 0 }
\end{lstlisting}

\paragraph{Charon's output.}
Charon adds a \texttt{RemainderByZero} UB check
(\texttt{var4 = (2 == 0)}, always false) and a subsequent \texttt{Assert}
(dropped by \texttt{charon2lean.py}).
The resulting \texttt{LLBCFunDef} body is:
copy \texttt{*self} $\to$ \texttt{var3}; compute \texttt{var2 = var3 \% 2};
compute \texttt{var0 = (var2 == 0)}; return.

\paragraph{Key challenge.}
Using \texttt{num-integer} as a \emph{dependency} (the natural way)
produces only opaque bodies — Charon cannot see into compiled
dependency crates. Running Charon on the crate source directly
resolves this, but requires access to the source (available in
the Cargo registry cache at \texttt{\textasciitilde/.cargo/registry/src/}).

\paragraph{Verification.}
Five theorems in \texttt{Translation/NumIntegerSpec.lean}:

\begin{lstlisting}[caption={NI1 (symbolic): Integer::is\_even<u32> verified.},
                   label={lst:nispec}]
theorem num_is_even_symbolic (n : Nat) :
    evalFun [] NumIsEvenU32Fun 10 [.uint n] at () |=
      .pureOutput (· = .bool_ (n % 2 == 0)) :=
  ⟨.bool_ (n % 2 == 0), (), rfl, rfl⟩
\end{lstlisting}

The proof is by \texttt{rfl}: the kernel evaluates
\texttt{n \% 2 == 0} symbolically (the UB check `2 == 0' reduces to
\texttt{false} for the literal divisor, leaving only the arithmetic).
Ground instances (NI2--NI4) and a nocrash theorem (NI5) complete the set.

\subsection{TypeScript Cross-Language Unification}
\label{sec:eval-ts}

The TypeScript module (\texttt{TypeScript/ElabSpec.lean}, 498 lines) now
covers 14 functions and includes 6 unification theorems:

\begin{itemize}[leftmargin=2em]
  \item \textbf{T1--T14}: pure functions (return42, id, neg, add, sub, abs,
        max, isZero, mul, min, not\_gate, clamp) and the looping
        \texttt{sum\_to} function proved correct at four ground instances
        (n = 0, 1, 5, 10).
  \item \textbf{U1--U6}: cross-language unification theorems asserting that
        both a Rust LLBC implementation and a TypeScript implementation satisfy
        the \emph{same} \ProgramSpec predicate.
        Notable: U5 pairs \texttt{sumToFun} (Rust LLBC loop) with
        \texttt{tsSumToFun} (TypeScript \texttt{while} loop), both
        proving \texttt{sum\_to(10) = 45} from the identical specification.
\end{itemize}

The TypeScript \texttt{while} loop required extending the \texttt{TSStmt}
AST with a single mutable-update constructor (\texttt{set\_}), added in
3 lines to \texttt{AST.lean} and \texttt{Elaborator.lean}.
All existing proofs continue to build without modification.

\subsection{Lines of Code}

\begin{table}[t]
\centering
\caption{Lines of code by category (21 source files, 5\,361 total lines).}
\label{tab:loc}
\small
\begin{tabular}{@{}lrr@{}}
\toprule
\textbf{Category} & \textbf{Files} & \textbf{Lines} \\
\midrule
Foundation (Types, Monad, Ownership)                      & 3  &  281 \\
Specification (ProgramSpec, Satisfies, Examples)          & 3  &  427 \\
Translation (LLBC, Elaborator, ElabSpec, Semantics, Adequacy) & 5 & 2{,}095 \\
Loop invariants (LoopInvariant, FactInvariant, FibInvariant)  & 3 &   748 \\
Charon pipeline (CharonDefs, CharonSpec)                  & 2  &  756 \\
Tactic automation (VerifyFun, Examples)                   & 2  &  186 \\
TypeScript (AST, Elaborator, ElabSpec)                    & 3  &  868 \\
\midrule
\textbf{Total (Lean)}                                     & \textbf{21} & \textbf{5\,361} \\
Charon pipeline script (\texttt{charon2lean.py})          & 1  &  461 \\
\bottomrule
\end{tabular}
\end{table}

\subsection{Proof-to-Code Ratios}

\begin{table}[t]
\centering
\caption{Proof-to-code ratios. ``Def lines'' counts definition bodies;
         ``Proof lines'' counts proof bodies; ``Ratio'' = Proof/Def.}
\label{tab:ratios}
\small
\begin{tabular}{@{}lrrr@{}}
\toprule
\textbf{Module} & \textbf{Def lines} & \textbf{Proof lines} & \textbf{Ratio} \\
\midrule
\texttt{Spec/Satisfies.lean} (combinators)      &  51 & 136 & 2.7$\times$ \\
\texttt{Translation/ElabSpec.lean} (LLBC pure)  &  89 & 428 & 4.8$\times$ \\
\texttt{Translation/CharonSpec.lean} (Charon)   &  40 & 290 & 7.3$\times$ \\
\texttt{Tactic/Examples.lean} (macro proofs)    &  22 &  16 & 0.7$\times$ \\
\texttt{Translation/LoopInvariant.lean}         &  63 & 291 & 4.6$\times$ \\
\texttt{Translation/FibInvariant.lean}          &  47 & 186 & 4.0$\times$ \\
\texttt{TypeScript/ElabSpec.lean}               & 110 & 387 & 3.5$\times$ \\
\bottomrule
\end{tabular}
\end{table}

The \texttt{Tactic/Examples.lean} ratio of $0.7{\times}$ reflects the
proof automation macros: 16 theorems proved in 16 lines (one per theorem)
against 22 lines of statements. Without the macros, the same 16 theorems
would require approximately 176 lines ($11{\times}$ overhead).

The \texttt{CharonSpec.lean} ratio ($7.3{\times}$) is higher than the
hand-crafted \texttt{ElabSpec.lean} ($4.8{\times}$) because Charon-generated
functions use more temporaries (copy locals, overflow-check booleans), making
each function body larger relative to its proof.

\subsection{Verification Build Time}

We measured \lakeCmd{lake build Theorems} on an Apple M-series laptop
with 16\,GB RAM:

\begin{itemize}[leftmargin=2em]
  \item Full build from scratch (with \lakeCmd{lake exe cache get}): $\approx 4$ minutes.
  \item Incremental build (only lean-pl-verify changed): $\approx 30$ seconds.
  \item Mathlib compilation from source (no cache): 45--60 minutes.
\end{itemize}

The bulk of the build time is Mathlib compilation; lean-pl-verify itself
takes approximately 25--30 seconds to elaborate. \texttt{FibInvariant.lean}
takes the longest single file ($\approx 8$ seconds) due to the inductive
proof over $n.\texttt{toNat}$ iterations and the \texttt{native\_decide} call.

\subsection{Reproducibility}

The complete artifact is reproduced by:

\begin{lstlisting}[language={},caption={Reproduction commands.},label={lst:repro}]
$ git clone <repository> lean-pl-verify
$ cd lean-pl-verify
$ lake exe cache get      # download Mathlib cache (~1 GB, first time only)
$ lake build Theorems     # verifies all 258 theorems, 0 sorry
\end{lstlisting}

Successful completion of \lakeCmd{lake build Theorems} with no errors and no
sorry warnings is the reproducibility certificate.

\subsection{Comparison with Related Work}

\begin{table}[t]
\centering
\caption{Comparison with Aeneas on key dimensions.}
\label{tab:comparison2}
\small
\begin{tabular}{@{}lll@{}}
\toprule
\textbf{Dimension} & \textbf{Aeneas~\cite{ho2022aeneas}} & \textbf{This work} \\
\midrule
Embedding strategy    & Shallow (generated Lean functions) & Deep (inductive AST) \\
TCB                   & Charon + Aeneas translator         & Lean kernel only \\
Proof style           & Proofs over translated functions   & Proofs over AST \\
Loop invariants       & Yes (manual + loop summaries)      & Yes (manual, fuel-parametric) \\
Overflow model        & Explicit (Primitives library)      & Explicit (IntBounds checks) \\
TypeScript support    & No                                 & Yes (33 theorems) \\
Cross-language unif.  & No                                 & Yes (6 theorems) \\
Automated pipeline    & Yes (Charon integration)           & Yes (\S\ref{sec:eval-charon}) \\
Relational semantics  & No                                 & Yes (Semantics + Adequacy) \\
Sorry count           & N/A                                & \textbf{0 of 258 theorems} \\
Unsafe Rust           & No                                 & No \\
\bottomrule
\end{tabular}
\end{table}

The two systems are complementary rather than competing. Aeneas provides a
production-grade pipeline with a large body of verified cryptographic
code~\cite{aeneascrypto2026}. Lean-pl-verify provides a smaller,
kernel-minimal foundation that supports cross-language reasoning, a
relational semantics with adequacy, and an end-to-end Charon pipeline
to the same proof framework.
Combining the Aeneas translation pipeline with the \ProgramSpec unification
framework would be a natural direction for future work.

%% ============================================================
%%  §10b  Agentic Development: AI-Assisted Formal Verification
%% ============================================================

\section{Agentic Development of the Verification Library}
\label{sec:agentic}

lean-pl-verify was developed with sustained assistance from Claude Code
(Anthropic's AI coding agent, powered by Claude Sonnet~4.6).
This section documents what the agent automated, what required human judgment,
and what this experience implies for AI-assisted formal verification.
We report this not as a claim that AI ``wrote'' the proofs, but as an honest
account of the human--AI division of labor in a non-trivial verification project.

\subsection{What the Agent Automated}

\paragraph{Boilerplate proof generation.}
The large majority of the 258 theorems follow one of four proof shapes
(§\ref{sec:rust-proofs}).  Once the first theorem of each shape was written
by hand and verified, the agent generated the remaining theorems in that
shape by pattern-matching on the function body and instantiating the template.
This was most valuable for:

\begin{itemize}[leftmargin=2em]
  \item The 48 \texttt{ElabSpec.lean} theorems: 16 functions × 2--4 properties each.
  \item The 47 \texttt{CharonSpec.lean} theorems: generated from the
        Charon-extracted defs using the \texttt{charon\_simp} tactic macro.
  \item The 33 \texttt{TypeScript/ElabSpec.lean} theorems: mirroring the
        Rust proofs for the TypeScript evaluator.
  \item The 6 \texttt{TypeScript/BugDetection.lean} theorems: bug-expose and
        fix-verify proofs generated by instantiating the same \rfl template.
\end{itemize}

\paragraph{Simp set construction.}
Identifying which simp lemmas are needed to reduce a given evaluator call is
tedious but mechanical.  The agent iterated on failing simp calls by reading
error messages, adding missing lemmas (\texttt{getElem?\_pos},
\texttt{List.getElem\_cons\_succ/zero}, \texttt{decide\_eq\_true/false}),
and verifying the fix.

\paragraph{The Charon pipeline (\texttt{charon2lean.py}).}
The 461-line Python translator from Charon JSON to Lean \texttt{LLBCFunDef}
terms was written almost entirely by the agent.
The agent handled the flat-list-to-CPS conversion, the AddChecked/SubChecked
desugaring, the overflow-check Assert elision, negative-integer literal
wrapping, and type deduplication — each discovered by reading the 104\,KB
LLBC JSON and adapting to Lean's parsing requirements.
Human involvement was limited to approving the overall approach and
unblocking specific Charon installation issues (the \texttt{rustc-dev}
component requirement).

\paragraph{Adequacy proof scaffolding.}
The 14-theorem adequacy proof in \texttt{Adequacy.lean} was largely
scaffolded by the agent after the human specified the top-level theorem
statement.  The agent proposed the induction structure, identified the
\texttt{split at h} and \texttt{change} patterns needed for do-notation
reductions, and handled the 10+ statement cases.

\subsection{What Required Human Judgment}

\paragraph{Architectural decisions.}
The core choices — shallow monadic embedding, fuel-based evaluation,
forward-backward pairs for \texttt{\&mut T}, the \texttt{ProgramSpec}
inductive, positional environments for TypeScript — were made by the human
author based on the research goals and prior PL verification experience.
These decisions cannot currently be automated because they require understanding
the paper's contribution claims and the mathematical elegance trade-offs.

\paragraph{Sorry elimination.}
The \texttt{prodRange\_factorial} and \texttt{factSafe\_implies\_ov}
sorry closures required human-directed proof strategy (right-peel lemma,
\texttt{interval\_cases} for factorial bounds).
The agent could fill in the tactic steps once the strategy was specified,
but could not discover the strategy independently.

\paragraph{TypeScript AST design.}
The decision to add \texttt{set\_} (mutable update) to \texttt{TSStmt}
rather than use de Bruijn rebinding tricks for loop variables required
understanding the proof implications of each choice.  The agent proposed
both options but deferred to the human on which produced cleaner proofs.

\paragraph{Paper structure.}
The high-level narrative, the choice of which theorems to highlight, and
the framing of the unification claim as a first-class contribution were
human decisions.

\subsection{Workflow and Persistent Memory Model}

\begin{figure}[h]
\centering
\small
\begin{tabular}{@{}ll@{}}
\toprule
\textbf{MEMORY.md entry type} & \textbf{Example} \\
\midrule
Proof pattern & \texttt{split at h} for nested match reductions \\
Known simp lemma & \texttt{List.getElem?\_cons\_zero} (not \texttt{getElem?\_pos}) \\
Architectural fact & FBPair encodes \texttt{\&mut T} via forward/backward pair \\
Build fix & \texttt{charon2lean.py}: op\_name can be dict \texttt{\{"Rem":"UB"\}} \\
Sorry status & \texttt{FactInvariant.lean}: 2 sorry → 0 sorry (resolved) \\
\bottomrule
\end{tabular}
\caption{Categories of entries in \texttt{MEMORY.md} (persistent cross-session state).}
\label{tab:memory}
\end{figure}

The development proceeded in a persistent conversation model: each session
began with the agent reading \texttt{MEMORY.md} (maintained via the
\texttt{Write}/\texttt{Edit} tools), which accumulated architectural facts,
proof patterns, discovered simp lemmas, and open task status.

\textbf{Error correction taxonomy.}
Table~\ref{tab:errors} categorises the types of errors the agent successfully
corrected versus where it needed human redirection.

\begin{table}[h]
\centering
\small
\caption{Agent error handling in lean-pl-verify development.}
\label{tab:errors}
\begin{tabular}{@{}lll@{}}
\toprule
\textbf{Error type} & \textbf{Outcome} & \textbf{Example} \\
\midrule
Missing simp lemma & Agent resolves & Added \texttt{getElem?\_cons\_succ} from error \\
Wrong fuel bound (off by 1) & Agent resolves & Increased fuel from 10 to 12 \\
Unicode in listings (UTF-8 byte) & Agent resolves & Added \texttt{lstdefinelanguage\{TypeScript\}} \\
Dict-form BinOp in Charon JSON & Agent resolves & \texttt{isinstance(op\_name, dict)} check \\
\midrule
Proof strategy (which lemma to use) & Human directs & Right-peel for \texttt{prodRange\_factorial} \\
Architectural choice & Human decides & \texttt{set\_} vs de~Bruijn rebinding for TS \\
Sorry closure strategy & Human directs & \texttt{interval\_cases} for factorial bounds \\
Context loss (re-proposing failed fix) & Human redirects & \texttt{MEMORY.md} partially mitigates \\
\bottomrule
\end{tabular}
\end{table}

Observed productivity gains:
\begin{itemize}[leftmargin=2em]
  \item \textbf{Build-debug cycles} were 3--5$\times$ faster: the agent
        read error messages, proposed fixes, and rebuilt without waiting for
        human interpretation of Lean's error output.
  \item \textbf{Repetitive proof patterns} (same proof shape × 16 functions)
        were handled in a single agent invocation rather than 16 manual edits.
  \item \textbf{Cross-file consistency} was maintained by the agent reading
        both the definition file and the proof file before suggesting changes.
  \item \textbf{Context-window management}: sessions compressed prior context
        automatically; \texttt{MEMORY.md} was the primary mechanism for
        preserving architectural decisions across compression boundaries.
        Structured proof-attempt logs (recording which approaches failed and why)
        would further improve multi-session continuity.
\end{itemize}

\subsection{Implications for Formal Verification}

The experiment suggests a productive division of labor:
humans specify the \emph{what} (theorem statements, proof strategies,
architectural choices) and agents handle the \emph{how} (tactic sequences,
simp sets, boilerplate generation, pipeline scripts).
This mirrors the relationship between a mathematician and a proof assistant
— the mathematician specifies the proof sketch; the assistant checks
the details — but with the agent absorbing an additional layer of mechanical
work above the kernel.

The main limitation observed was \textbf{context loss}: in long sessions,
the agent occasionally re-proposed solutions already known to fail,
requiring the human to redirect.
The persistent memory file (\texttt{MEMORY.md}) mitigated but did not
eliminate this.  Structured proof-attempt logs (recording which approaches
failed and why) would be a valuable engineering improvement.

We anticipate that agentic formal verification will become standard practice
as LLMs improve at understanding proof error messages and tracking proof
state across files.  lean-pl-verify demonstrates that the current generation
of agents can already handle a non-trivial 258-theorem library with
\textbf{0 sorry}.

%% ============================================================
%%  §11  Related Work
%% ============================================================

\section{Related Work}
\label{sec:related}

\subsection{Rust Verification in Lean 4}

\textbf{Aeneas}~\cite{ho2022aeneas} is the primary related system. It translates
safe Rust programs via Charon into pure Lean~4 (and F$^\star$, Coq, HOL4)
using a shallow monadic embedding. Aeneas and lean-pl-verify share the goal
of Lean~4-based Rust verification and the use of a state-error monad, but
they differ fundamentally in embedding strategy. In Aeneas, the Lean function
generated for a Rust function \emph{is} the functional model; its correctness
depends on the soundness of the Charon+Aeneas pipeline. In lean-pl-verify,
the Lean definition of \texttt{evalStmtFuel} \emph{is} the semantics; there
is nothing to be sound with respect to, beyond what the kernel computes.

Recent applications of Aeneas include verified cryptographic
primitives~\cite{aeneascrypto2026}: the BLAKE3, SHA-2, and Kyber
implementations in the HACL* library. These verify real-world, production
code at a scale that lean-pl-verify does not currently match. The advantage
of Aeneas in this context is the Charon pipeline's ability to extract from
arbitrary crates; our manual transcription limits scale.

\textbf{LOOM}~\cite{loom2026} develops a framework for multi-modal
foundational verification using Dijkstra monads in Lean~4. LOOM's goals
overlap with ours (language-agnostic specifications, Lean~4 kernel as trust
anchor), but LOOM operates on shallow program representations rather than
deep AST embeddings. The Dijkstra monad index provides automated WP
generation, which LOOM exploits for more automated proofs. Whether LOOM's
approach can be extended to cross-language unification theorems similar to
ours is an interesting open question.

\subsection{Rust Verification via SMT}

\textbf{Verus}~\cite{verus2023} is currently the most practical tool for
verifying real-world Rust programs, with successful applications to distributed
systems and operating system code. Verus embeds specifications directly into
Rust source code using a ghost-type discipline; the Z3 SMT solver handles
proof obligations. The principal limitation is the ``SMT cliff'': complex
proofs, especially those involving non-linear arithmetic or intricate data
structure invariants, can cause Z3 to time out. Verus does not support
TypeScript and does not produce Lean kernel-checked proofs.

\textbf{Creusot}~\cite{creusot2022} uses Why3 as a deductive verification
backend. Programs are annotated with Rust attributes, and Creusot generates
Why3 goals discharged by a portfolio of SMT solvers. It supports a wider
fragment of Rust than Verus (notably, trait objects and iterators), but
shares the reliance on SMT solvers in the TCB.

\subsection{Bounded Verification}

\textbf{Kani}~\cite{kani2022} is a bounded model checker for Rust built on
CBMC. It handles unsafe Rust and concurrency by bounding the execution depth
and encoding safety properties as assertions for a SAT/SMT solver. Kani
excels at finding concrete bugs (memory unsafety, integer overflow in specific
inputs) but cannot prove properties for all inputs or for loops with
unbounded iteration counts. It does not produce formal proof certificates.

\textbf{SEABMC} and related tools apply bounded symbolic execution to
Rust via LLVM IR. They share Kani's strengths and limitations regarding
unbounded verification.

\subsection{Separation Logic and Unsafe Rust}

\textbf{RustBelt}~\cite{rustbelt2018} is the foundational semantic model
for Rust's type system, developed in Coq using the Iris separation logic
framework~\cite{iris2015}. RustBelt proves that the Rust type system
enforces memory safety even in the presence of unsafe code, by providing
a semantic model for lifetimes and type invariants. The proofs are highly
manual and cover only a small fragment of the standard library.

Our work shares RustBelt's interest in formal semantics for Rust but
targets a complementary level: functional correctness of safe programs
rather than type-soundness of unsafe abstractions. Combining a
RustBelt-style unsafe-code foundation with our safe-code verification
layer would be a natural long-term goal, following the hybrid architecture
proposed by Jung~\etal~\cite{rustbelt2018}.

\subsection{Cross-Language and Multi-Language Verification}

Cross-language correctness has been studied in the context of compilation
verification~\cite{compcert2009}: if a compiler is verified to preserve
semantics, then properties proved about the source program hold for the
target. Our approach is different: we do not verify a compiler; we instead
provide a common semantic target into which multiple source languages are
independently embedded.

\textbf{Bedrock2} and related work in the certified programming languages
community has explored ``push-button'' verification of C code compiled to
machine code, establishing end-to-end correctness from source to binary.
The multi-language direction in lean-pl-verify is orthogonal: rather than
going deeper (closer to hardware), we go wider (across languages).

\textbf{F$^\star$}~\cite{fstar2016} supports multiple languages via a
common F$^\star$ extraction target. The Low$^\star$ sublanguage is used
to verify C programs; KReMLin extracts to C. This is structurally similar
to our approach (common monad, multiple source embeddings) but F$^\star$'s
multi-language support is directed at compilation rather than cross-language
property sharing.

\subsection{Proof Automation for Program Verification}

Lean~4's \texttt{grind} tactic~\cite{lean4} combines congruence closure,
E-matching, and linear arithmetic; it is designed precisely for program
verification goals. We deliberately avoided \texttt{grind} in lean-pl-verify,
preferring explicit \texttt{simp only} calls that make the proof state
predictable. Whether \texttt{grind} could further compress the loop invariant
proofs is unexplored.

The Lean Squad project~\cite{leansquad2026} reports near-zero human labour
verification of Rust programs using an LLM agent that drives the Lean LSP
via the Model Context Protocol. This is architecturally different from our
work (agent vs.\ manual proof writing) but shares the goal of using Lean~4
as the verification target. The agent uses Aeneas for the Rust-to-Lean
translation; combining the agent with our deep embedding would be an
interesting experiment.

\subsection{Domain-Specific Languages for Verification}

\textbf{Dafny}~\cite{dafny2010} provides an integrated verification language
where specifications (\texttt{ensures}/\texttt{requires}/\texttt{invariant})
are interleaved with program code. It supports loops, recursive functions,
and SMT-backed automation. Dafny does not target Rust or TypeScript.

\textbf{DSLean}~\cite{dslean2026} develops a framework for type-correct
interoperability between Lean~4 and external DSLs, using Lean's elaboration
and metaprogramming to enforce typing in real time. This is related to our
macro-based proof automation (\S\ref{sec:automation}) but aimed at DSL
interoperability rather than proof compression.

%% ============================================================
%%  §12  Conclusion
%% ============================================================

\section{Conclusion}
\label{sec:conclusion}

\subsection{Summary}

We have presented lean-pl-verify, a Lean~4 library for cross-language formal
verification. The central contribution is a pair of language-agnostic components
— the \RustM execution monad and the \ProgramSpec specification language — into
which programs from different source languages are independently embedded. Once
in this common form, all specification and proof tools apply uniformly.

The library contains \textbf{258 kernel-checked theorems} across 13 modules,
with \textbf{0 sorry}:
48 LLBC Rust theorems covering 16 functions (hand-crafted LLBC);
14 theorems for the sumTo loop invariant;
14 for factorial; 9 for Fibonacci;
18 for full adequacy (soundness + completeness) of the relational semantics;
57 over Charon-extracted function definitions (18 functions including \texttt{pow} and \texttt{gcd});
5 for \texttt{Integer::is\_even<u32>} from the published \texttt{num-integer} v0.1.45 crate;
39 TypeScript theorems including 6 cross-language unification results
and a 6-theorem bug-detection case study;
17 ProgramSpec combinators; 15 monad-level examples; 4 ownership laws;
2 monad laws; and 16 proof automation demos.
Full adequacy (soundness and completeness) is proved with 0 \sorry.
The only trust anchor for proof correctness is the Lean~4 kernel;
\lakeCmd{lake build Theorems} is the complete reproducibility certificate.

\subsection{What Was Harder Than Expected}

Three aspects of the development were more difficult than anticipated.

\textbf{Fuel arithmetic.} Counting the correct fuel bound for a multi-step
loop body requires tracking how fuel is consumed at each interpreter level.
Off-by-one errors in this accounting were the most common source of proof
failures. The \texttt{hpre := rfl} technique (\S\ref{sec:rust-proofs})
emerged to isolate the prefix and loop phases.

\textbf{Macro parsing in Lean~4.} The two-consecutive-term parsing issue
(\S\ref{sec:automation}) was genuinely surprising. Lean~4's parser is a
recursive descent parser with Pratt-style operator precedence; the interaction
between macro syntax rules and term-level parsers is not always intuitive.

\textbf{Adequacy proof patterns.} The \texttt{split at h} tactic for nested
\texttt{match} reductions, and the \texttt{change} tactic for do-notation
reduction steps, were non-obvious. Both are documented in
\S\ref{sec:semantics} as reusable patterns for future adequacy proofs.

\subsection{Limitations}

\textbf{Safe Rust only.} Raw pointers, \texttt{unsafe} blocks, and
concurrency are outside the model. For unsafe code, a separation logic
layer (Iris~\cite{iris2015} or a Lean~4 analogue) would be needed.

\textbf{Charon pipeline scope.} The Charon pipeline is demonstrated on
18 functions, including non-trivial algorithms (\texttt{pow}, \texttt{gcd})
that exercise Charon's UB-check infrastructure (\texttt{Rem} desugaring,
overflow guards). Scaling to real-world crates with hundreds of
functions, trait objects, and lifetimes requires further engineering.

\textbf{Loop invariants are manual.} The inductive invariant patterns for
\texttt{sumTo}, \texttt{fact}, and \texttt{fib} are provided by the proof
engineer. Automated invariant synthesis for kernel-verified systems is an
open problem.

\textbf{Completeness of the TypeScript semantics.} A relational semantics
for the TypeScript AST (analogous to the LLBC \texttt{EvalStmt} relation)
and a full adequacy proof for the TypeScript evaluator are natural next steps,
following the pattern of \texttt{Translation/Adequacy.lean}.

\subsection{Future Work}

\textbf{Scaling to real-world Rust crates.} We verified \texttt{Integer::is\_even<u32>}
from \texttt{num-integer} v0.1.45 by running Charon directly on the crate source.
The remaining \texttt{num-integer} functions use Stein's binary GCD, requiring
\texttt{trailing\_zeros} intrinsics and signed bitwise AND, not yet in our
evaluator's scope. Three concrete steps would enable broader coverage:
(i)~adding signed bitwise operations (\texttt{BitAnd/BitOr/BitXor} for \texttt{Int});
(ii)~modelling compiler intrinsics (\texttt{trailing\_zeros}, \texttt{count\_ones})
as axiomatised functions with proved properties; and
(iii)~automating loop invariant generation via an LLM agent loop —
the agent writes a candidate invariant, the kernel rejects or accepts, and the
loop iterates. This last step is the highest-leverage direction: the manual
invariant work is the dominant cost in verifying loop-heavy code.

\textbf{More languages.} The \ProgramSpec framework applies to any language
whose semantics can be embedded as a \RustM computation. Python, C, or WASM
would be natural targets; each would contribute new cross-language unification
theorems.

\textbf{Unsafe Rust.} Combining this work with an Iris-based model for
unsafe abstractions (following the hybrid architecture of~\cite{rustbelt2018})
would extend the verified fragment to programs using \texttt{Vec},
\texttt{Arc}, and similar standard library types.

\textbf{Automated loop invariants.} Integrating an invariant synthesis
loop — perhaps using an AI agent in the style of the agentic development
described in \S\ref{sec:agentic} — with the kernel-checked proof structure
would reduce manual effort for loop-heavy programs.

\textbf{Performance verification.} Extending \ProgramSpec with fuel-based
complexity specifications — ``this function terminates in at most $f(n)$
steps'' — would enable worst-case complexity proofs for verified programs.

\subsection{Final Remark}

Formal verification of systems software is hard. Tools exist for Rust, for C,
for functional languages — but they do not talk to each other, and proving
that a TypeScript function is correct in the same way as a Rust function
means starting from scratch. This paper argues that the gap is smaller than
it looks: a single monad and a single specification language, carefully
designed, can span both.

The \textbf{258 kernel-checked theorems} in lean-pl-verify are evidence for
this claim. The Charon pipeline shows the framework connects to real Rust
tooling. The relational semantics shows the interpreter is not an arbitrary
artifact but has a principled denotation. And the six cross-language
unification theorems are the clearest expression of the core idea: the
specification does not know or care which language the program came from.


\bibliographystyle{ACM-Reference-Format}
\bibliography{references}

\end{document}
